\documentclass{article}
\PassOptionsToPackage{numbers,sort&compress}{natbib}
\usepackage[dblblindworkshop]{neurips_2026}
\usepackage[T1]{fontenc}
\usepackage[utf8]{inputenc}
\usepackage{amsmath,amssymb}
\usepackage{graphicx}
\usepackage{booktabs,tabularx}
\usepackage{algorithm,algpseudocode}
\usepackage{microtype}
\usepackage{hyperref}
\hypersetup{hidelinks,pdfauthor={},pdftitle={LEAF: Does Agentic Search Earn Its Oracle Budget?}}
\title{Do LLM Proposers Improve Budgeted Molecular Search?}
\workshoptitle{Workshop title to be confirmed}
\author{Anonymous Author(s)}
\date{}

\begin{document}
\maketitle

\begin{abstract}
Does LLM-driven evolutionary proposal improve molecular search when oracle-query budgets and acquisition functions are matched? LEAF tests this question on a 30,044-candidate ionizable-lipid space using a frozen Gaussian-process oracle fitted to real AGILE potency labels and a separately updated active-learning surrogate. Four language models, each evaluated over eight seeds and three splits, show no significant advantage over genetic-algorithm or random-search baselines. A single-objective reanalysis restricted to chemotype-cold's one validated readout (dropping a failing secondary assay) confirms the deficit persists but at a far smaller effect size than the two-objective number suggested. Small local models produce predominantly duplicate proposals; direct instrumentation of the acquisition-selected, budget-charging queries (not just the raw proposal pool) shows this is worse than the pool-level rate suggests -- 79\% of one model's real oracle queries exactly duplicate an already-evaluated candidate, meaning acquisition concentrates on rather than filters out duplicates. A 14B model removes this failure mode but shows no evidence of either a difference from, or equivalence to, the baselines under a pre-specified effect-size threshold. Gemini Flash-Lite underperforms both baselines on all three splits. A diversity-aware prompt mitigates duplication without establishing an advantage. A matched de-duplication intervention, re-run at full statistical power ($n=8$) on the random split after external review, moves both baseline comparisons from significant to non-significant without reversing either median, at the cost of requiring far more proposal rounds to spend the same budget -- and, for two of eight seeds, being unable to spend the full budget at all. Extending this check (exploratory, $n=4$) to the other two splits shows the effect is not uniform: scaffold-cold narrows sharply, matching the random split's direction, while chemotype-cold's deficit barely moves, suggesting duplication is not the dominant mechanism on that split. A stricter de-duplication criterion that excludes all previously evaluated candidates, not only the current archive (exploratory, $n=4$), flips the median difference against static-GA from negative to positive and substantially narrows the deficit against random search, though wide, zero-crossing CIs at this sample size mean this is directional evidence, not a significance claim. A stricter de-duplication criterion excluding all previously evaluated candidates (not only the current archive) and a richer prompt exposing real objective scores each independently flip the median deficit against static-GA positive at $n=4$ -- directional, not significant, and the two checks share a convergent, non-independent seed we verified and disclose rather than let inflate the apparent evidence. A follow-up leakage probe finds the same base model has no zero-shot discriminative signal about real potency at all (a degenerate constant output, not merely a low correlation), weighing against a memorization explanation; a rank-only variant of the richer prompt matches or exceeds the raw-scores result, suggesting relative preference rather than magnitude-anchoring drives the effect. A perturbation check that randomly scrambles the score-to-fragment mapping shown in the richer prompt produces a statistically indistinguishable aggregate result, but inspecting real per-query trajectories shows this reflects the search space's narrow attainable frontier (different paths reconverging), not a uniform failure to use the mapping -- one seed shows a genuine round-by-round null effect, three show real behavioral divergence between conditions. A standalone evidence-conflict probe with a deliberately falsified quality ranking finds the model's proposals follow neither the false label (14.8\%, no different from chance) nor the real quality ranking (0\%), but instead cluster on whichever archive member is listed first (74\%), a sharper and more specific finding than either reading. A genuine reasoning-tuned model (native extended inference-time reasoning, not a prompted rationale) engages substantively with real injected scores when given them (81.6\% of its reasoning traces mention the objective by name, versus 0\% under an uninformative default prompt) yet still does not outperform a smaller non-reasoning model given the identical scored prompt, at 5--10$\times$ the latency. A first exploratory test of the paper's EHVI acquisition path (fixed after a dispatch bug was found) shows a similar deficit to UCB, not a reversal. The Verifier ran in every experiment but was non-contributory (no real toxicity assay yields zero-width conformal intervals); the Hypothesis Ledger was enabled in the primary sweep, where it remained empty for structural reasons, and was deliberately disabled in later isolation experiments. Zero-shot potency scoring with Qwen3-8B is uncorrelated with experimental labels, unlike the fitted oracle. These model-on-model results do not establish biological efficacy or a universal limitation of agents. They provide a budget-matched evaluation protocol, negative evidence across tested configurations, and an audit showing how ordinary implementation errors can manufacture apparent agentic gains.
\end{abstract}

\section{Introduction and related work}
\label{sec:introduction}
Ionizable lipids strongly influence lipid nanoparticle (LNP) delivery efficiency, and combinatorial changes to head groups, linkers, and tails create large design spaces from established synthetic routes \citep{xu2023ugi}. LEAF asks a methodological question rather than claiming improved lipid therapeutics: \emph{does an LLM evolutionary proposer improve the Pareto frontier over conventional search when the oracle, acquisition function, and query budget are identical?} Cheap, well-characterized combinatorial objectives offer little long-horizon reasoning or heterogeneous evidence for an LLM to exploit. Sparse labels, literature-grounded objectives, and multi-step synthesis may be more favorable settings, but are not tested here.

\textbf{Lipid design.} Combinatorial Ugi three-component-reaction (Ugi-3CR) libraries and AGILE's supervised screening pipeline precede agentic approaches \citep{xu2023ugi,xu2024agile}; synthesizable-lipid generative modeling provides another alternative \citep{zhu2024generative}. LipoAgent \citep{li2025lipoagent} coordinates fine-tuned TxGemma and Qwen Predict/Verify agents for library screening. Its published method includes toxicity-conditioned efficiency prediction; the efficiency-only code path inspected in the original LEAF record cannot establish that the full system lacks safety modeling. Its task still differs from matched-budget search. LUMI-lab combines foundation models with robotic active learning, screens over 1,700 LNPs in vitro, and validates candidates in vivo \citep{lumilab2026}. Its LUMI-6 result (20.3\% lung-epithelium gene editing after intratracheal administration) and brominated-tail/potency association do not establish a structure-to-organ-tropism rule.

\textbf{Agents and optimization.} The Virtual Lab, ProtAgents, protein-design swarms, and AI co-scientist exemplify collaborative biomolecular reasoning \citep{swanson2025virtuallab,protagents2025,swarmdesign2025,gottweis2026}. ChemCrow and Coscientist established chemistry-tool use \citep{bran2023chemcrow,boiko2023coscientist}; ToolMol motivates tool-scored crossover/mutation \citep{toolmol2026}. LEAF combines this proposal pattern with Gaussian-process (GP) active learning and a continuous Tanimoto kernel \citep{gessner2024,srinivas2009,rasmussen2006,settles2012}. Unlike library-only retrospective evaluation, generated tuples require a generalizing evaluation oracle, distinct from the acquisition surrogate. Batch EHVI, conditional generative multi-objective Bayesian optimization, and LLM-derived evolutionary operators inside a multi-objective evolutionary algorithm are complementary alternatives that this paper does not evaluate; we do not cite specific papers for these here, having been unable to independently verify author/venue details for candidate references before this submission.

\textbf{Contributions.} We specify a matched-budget test with an explicitly pre-committed negative-result interpretation; evaluate four fully replicated LLM configurations, a diversity prompt, a non-LLM bandit, and exploratory interventions; diagnose duplicate-driven waste and, separately, the non-contributory Verifier and the empty-then-deliberately-disabled Hypothesis Ledger; and document the correction of false-positive-producing implementation errors. Although the architecture includes a debate-style Verifier and a literature-grounded Hypothesis Agent, the Verifier was functionally non-contributory under the available data, and the Hypothesis Ledger produced no validated entries in the primary sweep. Interpretability is a possible qualitative benefit, not an independently established experimental gain.

\section{Method and experimental protocol}
\label{sec:method}
\subsection{Search space, objectives, and separate oracles}
The general lipid representation is $x=(h,\ell,t_1,t_2,u,b)\in\mathcal X$, indexing head group, linker, tails, unsaturation, and halogenation. The implemented Ugi-3CR assembler instead uses three reagent slots $(A,B,C)$: amine, aldehyde, and isocyanide. Molecular embeddings are Morgan fingerprints. The general optimization specification is
\begin{equation}
 f(x)=\bigl(f_{\mathrm{pot}}(x),-f_{\mathrm{tox}}(x)\bigr),\qquad
 f_{\mathrm{aug}}(x)=\bigl(f_{\mathrm{pot}}(x),-f_{\mathrm{tox}}(x),f_{\mathrm{syn}}(x)\bigr).
 \label{eq:objectives}
\end{equation}
The synthesizability axis is conditional on non-negligible variance across $\mathcal X$; a nearly constant route-validity score cannot meaningfully enlarge a Pareto frontier. It is never instantiated in any reported run (confirmed by inspection of every experiment configuration): $f_{\mathrm{aug}}$ is a specified but unused generalization of Eq.~\ref{eq:objectives}, not a description of the actual optimization target. Tropism heuristics are secondary filters, not optimization axes or evidence of organ selectivity. Table~\ref{tab:objectives} gives the exact, verified objective specification for every reported comparison, rather than deferring it to configuration files.

\begin{table}[t]
\caption{Exact objective specification for every reported experiment family. All entries verified directly against \texttt{leaf/orchestrator/loop.py} (\texttt{LoopConfig} defaults) and every \texttt{leaf/configs/*.yaml} file used to produce a reported result; no experiment overrides these defaults.}
\label{tab:objectives}
\centering
\begin{tabular}{@{}llllll@{}}
\toprule
Experiment family & Objectives & Direction & Normalization & Reference point & Toxicity/synth. \\
\midrule
All reported runs & \texttt{potency}, & maximize & GP internal & $(-5.0,-5.0)$, & Toxicophore rule \\
(primary sweep, GA, & \texttt{potency\_raw} & both & \texttt{normalize\_y} & fixed identically & used only to \\
random search, bandit, & (real AGILE & (both & (per-head & across every & screen the \\
diversity prompt, & \texttt{expt\_Hela}/ & signs & $z$-score, fit & split/method/seed & Verifier pool, \\
de-duplication check, & \texttt{expt\_Raw} & $=+1.0$) & on that split's & before any run) & never a third \\
Groq/Gemini) & labels) & & training data) & & objective axis \\
\bottomrule
\end{tabular}
\end{table}

Both objectives are real, correlated transfection-potency readouts from different assays (HeLa and RAW264.7), not independent efficacy axes; every reported hypervolume is therefore a two-objective potency-only quantity, not a potency/toxicity/synthesizability trade-off, and no toxicity or synthesizability signal ever enters the optimized vector in any experiment in this paper. The reference point is identical across every split within a given budget configuration, so within-budget hypervolume values are comparable across splits and methods; cross-budget comparisons (e.g.\ $B=400$ vs.\ $B=40$ preliminary conditions) are not, since a larger budget can access more of the same fixed frontier.

The frozen evaluation oracle $\mathcal O_{\mathrm{GT}}$ represents the laboratory in a retrospective simulation. It must generalize to assembled molecules absent from the labeled library. For split-specific experiments it is fitted only to that split's training partition and then frozen; a full-data fit must not erase the split distinction. The active-learning surrogate $\mathcal O_{\mathrm{AL}}$, using the same GP architecture, is refitted each round on acquired labels only. Its posterior drives acquisition and, where applicable, attribution and verification. Querying a surrogate for its own training targets would yield no new information; the two objects are therefore explicitly separated. The evaluation remains model-on-model, not experimental validation.

\begin{figure}[t]
\centering
\includegraphics[width=\linewidth]{Figure 1.pdf}
\caption{LEAF's closed-loop architecture. The Proposer (one of four interchangeable methods: LLM, static-GA, random, or Thompson bandit) generates a pool of $m=10$ candidates over the real fragment slots (A, B, C); the active-learning surrogate $\mathcal O_{\mathrm{AL}}$ (GP, Tanimoto kernel) scores the pool at no oracle-query cost; the Verifier applies the deterministic RDKit toxicophore rule; UCB ($\beta=2$) acquisition selects $|C_t|\le5$ candidates, the only step (gold arrows) that charges the shared $B=40$ query budget against the frozen oracle $\mathcal O_{\mathrm{GT}}$ (fit on the AGILE training split, real HeLa/RAW264.7 assay heads); the archive and surrogate then update and refit. The Hypothesis Agent's counterfactual tests (also gold, also budget-charging) periodically update the Hypothesis Ledger, which can feed back into the Proposer. This rendering accurately reflects Algorithm~\ref{alg:leaf}'s budget accounting: no generic drug-discovery iconography or mislabeled free/charged steps.}
\label{fig:architecture}
\end{figure}

\subsection{Proposal, acquisition, and secondary agents}
At round $t$, the LLM receives the Pareto archive $A_{t-1}$, real fragment vocabulary, and Hypothesis Ledger $H_{t-1}$, and proposes a pool $P_t$ using crossover and single-fragment mutation. A schema-constrained call returns one JSON candidate; reagent values are checked character-for-character against the vocabulary before assembly. Rejection/backfill rates distinguish genuine model output from fallback. The surrogate scores the pool without oracle-query cost. After verification removes abstentions, scalarized multi-objective UCB selects $q$ candidates using the configured exploration coefficient $\beta=2.0$. Only selected candidates incur evaluation-oracle charges. Static-GA and random search use the identical acquisition function and $q$ -- ``random search'' therefore names the \emph{proposal} operator (uniform vocabulary draws), not the full baseline: candidates it proposes are still budget-charged only after surviving the same UCB (or EHVI) selection as every other method, so a more literal name would be ``random-proposal + UCB.'' We retain ``random search'' as the conventional term for this baseline in the optimization literature, but flag this precisely rather than let the name imply a blind, unselected sampling procedure. The general design permits $m=|P_t|\gg q$; the primary LLM experiment uses $m=10$, $q=5$. EHVI is implemented, but a dispatch bug meant all reported historical experiments used UCB; no full EHVI robustness sweep is claimed. Exact per-method operator, backfill, and duplicate-handling detail is in Appendix~\ref{app:baselines}; exact LLM prompts, sampling parameters, and model identifiers are in Appendix~\ref{app:llm-repro}.

\textbf{Reader's note (external review).} This paper's central, evaluated contribution is the proposer comparison in \S3. The Verifier was enabled (\texttt{use\_verifier: true}) in every experiment in this paper, without exception, but is functionally non-contributory throughout: AGILE contains no experimental toxicity labels, so the conformal calibration set satisfies \texttt{cal\_scores = cal\_labels} exactly, producing zero-width intervals -- the Verifier reduces to a deterministic toxicophore filter and contributes no additional epistemic signal, in every run, not only some. The Hypothesis Ledger was enabled in the primary sweep and early ablations, where it remained empty across all 72 runs for a structural reason (\S3.4), and was deliberately disabled (\texttt{use\_hypothesis: false}), not silently inactive, in every later isolation experiment (de-duplication, EHVI, the single-assay reanalysis, and the enriched-prompt checks) specifically to focus evaluation on the proposer and acquisition function. Neither finding establishes that debate-based verification or hypothesis feedback offer no benefit in a configuration where a real toxicity signal exists or counterfactual matching is not Pareto-restricted -- these are implementation-coverage findings about what was tested, not general claims about the components. Readers focused on the primary proposer-vs-baseline question can skip to \S3 without loss of continuity.

The proposed Verifier aggregates independently prompted toxicologists from at least two base models using split conformal prediction \citep{angelopoulos2023}. Intervals straddling the accept/reject boundary cause abstention. Adaptive sampling breaks ordinary exchangeability; weighted/Mondrian recalibration is performed each round and coverage is tracked empirically, not assumed guaranteed. Agreement above 95\% for same-model instances flags inadequate diversity. In the real runs, however, no AGILE toxicity assay supports a learned head: \texttt{cal\_scores = cal\_labels} reproduces the same deterministic RDKit rule. Intervals have zero width, and repeated rule outputs are not independent epistemic evidence.

Every $k$ rounds the Hypothesis Agent uses literature passages and Integrated Gradients/KernelSHAP attributions \citep{sundararajan2017,lundberg2017} to propose structure-property rules. Validation requires matched counterfactual pairs evaluated by the frozen oracle, a minimum pair count, and a paired Wilcoxon test; querying the acquisition surrogate would merely recover its own fit. All pair queries count toward $B$. A mandatory bare-LLM leakage probe conditions any claimed rediscovery of LUMI-lab's brominated-tail/potency association: detected prior knowledge permits only a consistency check. The real vocabulary has no halogenated reagents, so that experiment was not run. Appendix~\ref{app:agents} specifies both secondary components.

\begin{algorithm}[t]
\caption{Budget-accounted LEAF round}
\label{alg:leaf}
\begin{algorithmic}[1]
\Require Archive $A_{t-1}$, ledger $H_{t-1}$, acquired data $D_{t-1}$, remaining budget $B_t$, frozen $\mathcal O_{\mathrm{GT}}$, surrogate $\mathcal O_{\mathrm{AL}}$
\State $P_t\gets\mathrm{Proposer}(A_{t-1},H_{t-1},\mathcal X)$
\State Optionally remove exact archive duplicates from $P_t$ (all methods)
\State $\hat f\gets\mathcal O_{\mathrm{AL}}.\mathrm{posterior}(P_t)$
\State $P_t'\gets\mathrm{VerifierSurvivors}(P_t,\hat f)$
\State $C_t\gets\mathrm{Acquisition}(P_t',\hat f,\min(q,B_t,|P_t'|))$
\State $D_t\gets D_{t-1}\cup\{(x,\mathcal O_{\mathrm{GT}}(x)):x\in C_t\}$; $B_t\gets B_t-|C_t|$
\State Refit $\mathcal O_{\mathrm{AL}}$ on $D_t$; $A_t\gets\mathrm{NonDominatedSort}(D_t)$
\State $H_t\gets H_{t-1}$
\If{$t\bmod k=0$ and budget permits the minimum complete pair set}
\State Construct matched counterfactual query set $Q_t$ with $|Q_t|\le B_t$
\State $y_t\gets\mathcal O_{\mathrm{GT}}.\mathrm{query}(Q_t)$; $B_t\gets B_t-|Q_t|$
\State $H_t\gets\mathrm{HypothesisTest}(A_t,H_{t-1},Q_t,y_t,\text{attributions},\text{literature})$
\EndIf
\State \Return $A_t,H_t,D_t,B_t$; repeat while budget remains
\end{algorithmic}
\end{algorithm}

\subsection{Data, comparisons, and pre-committed hypotheses}
The real AGILE labeled set comprises 1,200 lipids: a complete $20\times12\times5$ factorial design. Inverting the assembler on AGILE's 12,276-candidate unlabeled virtual library recovers all reagent triples and adds 9/62/9 amine/aldehyde/isocyanide values, yielding $29\times74\times14=30{,}044$ valid combinations. Decomposition reproduces all 1,200 known molecules exactly; full-space RDKit round-trip validation succeeds throughout. Only the original 1,200 molecules are labeled, so the remaining 28,844 require oracle extrapolation. At $B=80$, $B/|\mathcal X|\simeq0.27\%$. Earlier experiments use the smaller labeled space to diagnose saturation. The broader data plan includes a roughly 1,080-lipid Ugi library, within-source rank normalization plus a source indicator for cross-assay harmonization, and checksum logging. An approximately 1,900-lipid LUMI-lab pilot dataset is reserved for conditional external validation if released; it is not development data or a completed experiment.

Random, scaffold-cold, and chemotype-cold splits test increasing extrapolation. The pre-specified validity gate is held-out Spearman $\rho\ge0.3$; a failing readout/split is uninterpretable. The primary LLM protocol uses $n=8$ paired seeds per split, $B=40$, $q=5$, $m=10$, default temperature 0.7, and an 80-call proposal workload. Baselines share the vocabulary, evaluation oracle, budget, and acquisition: static-GA, random search, and a non-evolutionary Thompson-sampling bandit fitting per-fragment Beta posteriors afresh from the Pareto archive. Ablations comprise proposer/acquisition alone, addition of the Verifier, full LEAF, disabled Ledger feedback, and a native binary filter replacing debate/conformal verification. LipoAgent is assessed only in a separate prediction task.

The primary statistic is the median paired hypervolume difference, with a 95\% bootstrap confidence interval (CI); two-sided Wilcoxon signed-rank tests are secondary. At $n=3$, the smallest exact two-sided $p$ is 0.25; at $n=4$, 0.125; and at $n=8$, 0.0078125. Thus small pilots cannot establish $p<0.05$, while eight seeds permit, but do not guarantee, adequate power. Reported tests are unadjusted. Other planned metrics are Recall@$k$ against known high-potency lipids (MC3/SM-102/ALC-0315-level benchmarks), hypothesis precision, round-wise conformal coverage, and inter-agent agreement; missing metrics are not implied to have been measured.

\textbf{Hypotheses and fallback.} H1 predicts higher hypervolume than both baselines, with a larger gain on chemotype-cold than random splits. H2 predicts improved sample efficiency from Ledger feedback; H3 predicts lower false acceptance at matched coverage from recalibrated debate versus a binary gate. H4 predicts leakage-conditional rediscovery of the brominated-tail/potency association. The recorded pre-commitment is to interpret an H1 null as no measured design-efficiency benefit, rather than reframe it as a positive lipid-design result. \textbf{Caveat (external review):} this pre-commitment was practiced during development (H1's fallback interpretation was written before the corresponding experiments ran, \S3.4's audit) but was not separately timestamped and archived as an external, independently verifiable artifact; readers should weigh ``pre-committed'' as an internal development discipline, not a dated, third-party-verifiable pre-registration.

\section{Results and mechanism analysis}
\label{sec:results}
\subsection{Oracle validity and zero-shot prediction}
The frozen oracle predicts held-out HeLa potency with Spearman correlations 0.68/0.66/0.49 on random/scaffold-cold/chemotype-cold splits, all above threshold (Figure~\ref{fig:performance}A). The RAW264.7 readout is weaker; its chemotype-cold correlation is 0.23 and fails the gate. AGILE supplies no toxicity assay. These validity results support controlled surrogate comparisons, not biological accuracy on every extrapolated molecule.

\textbf{HeLa-only reanalysis of chemotype-cold (external review item \#1).} Because RAW264.7 fails the validity gate specifically on chemotype-cold, the two-head hypervolume reported for that split throughout \S3.2 is not fully interpretable on its failing axis. We re-ran the identical primary protocol ($n=8$, $B=40$, $q=5$, $m=10$, Llama 3.1 8B) restricted to the single validated readout (\texttt{potency}/HeLa only, single-objective hypervolume): the deficit persists but shrinks sharply in absolute terms -- median $-3.93$ against \emph{both} static-GA and random search (identical baseline hypervolumes across all eight seeds under this restricted single-objective metric), Wilcoxon $p=0.016$ against both. This is the principal, fully-interpretable chemotype-cold result; the original two-head number ($-31.2$/$-46.3$, \S3.2) remains diagnostic but should not be read as the split's primary validated finding. The direction and significance survive removing the uninterpretable axis entirely, which strengthens rather than undermines the qualitative conclusion (the deficit is not an artifact of the failing RAW264.7 head), but the true effect size on chemotype-cold is far smaller than the two-head number suggested -- roughly a 28\% relative deficit against a baseline hypervolume of $\approx14$, not the larger-looking absolute gap the two-head numbers imply.

In a separate scoring test, Qwen3-8B-a base model used by LipoAgent-assigns zero-shot efficiency scores from 1 to 10 at temperature zero, with no label access. On 50 real held-out random-split lipids, it has 0\% rejection and essentially no potency correlation ($\rho=0.014$, $p=0.92$), compared with the reported fitted-oracle correlation of 0.68 on the matched molecules. This is a base-model zero-shot control, not a replication of LipoAgent's fine-tuning, toxicity-conditioned objective, exact prompts, TxGemma component, or ensemble. It therefore cannot refute the published system's predictive performance.

\subsection{Matched-budget proposal comparisons}
\begin{table}[t]
\caption{Complete eight-seed comparisons. Each cell gives median paired hypervolume difference (method minus baseline), with the two-sided Wilcoxon $p$ beneath it. GA: static-GA; RS: random search. Diversity-prompt exact $p$-values (previously reported only as ``n.s.'') are now included, computed directly from the same real per-seed results (Appendix~\ref{app:tables}).}
\label{tab:summary}
\centering
\begin{tabular}{@{}llccc@{}}
\toprule
Method & Baseline & Random & Scaffold-cold & Chemotype-cold \\
\midrule
Llama 3.1 8B & GA & \shortstack{$-47.1$\\$0.055$} & \shortstack{$-24.8$\\$0.078$} & \shortstack{$-31.2$\\$0.039$} \\
 & RS & \shortstack{$-42.0$\\$0.0078$} & \shortstack{$-28.8$\\$0.039$} & \shortstack{$-46.3$\\$0.016$} \\
Qwen2.5 7B & GA & \shortstack{$-40.9$\\$0.078$} & \shortstack{$-21.6$\\$0.109$} & \shortstack{$-29.2$\\$0.016$} \\
 & RS & \shortstack{$-29.1$\\$0.016$} & \shortstack{$-34.3$\\$0.195$} & \shortstack{$-43.4$\\$0.023$} \\
Qwen2.5 14B & GA & \shortstack{$+1.8$\\$0.742$} & \shortstack{$+18.3$\\$0.313$} & \shortstack{$-55.1$\\$0.0078$} \\
 & RS & \shortstack{$-4.3$\\$0.742$} & \shortstack{$+33.9$\\$0.250$} & \shortstack{$-77.0$\\$0.0078$} \\
Gemini Flash-Lite & GA & \shortstack{$-91.8$\\$0.0078$} & \shortstack{$-66.6$\\$0.0078$} & \shortstack{$-24.8$\\$0.0078$} \\
 & RS & \shortstack{$-85.5$\\$0.0078$} & \shortstack{$-71.1$\\$0.0078$} & \shortstack{$-60.2$\\$0.0078$} \\
Diversity prompt & GA & \shortstack{$-32.5$\\$0.078$} & \shortstack{$-15.6$\\$0.148$} & \shortstack{$-10.1$\\$0.25$} \\
 & RS & \shortstack{$-25.4$\\$0.016$} & \shortstack{$-18.0$\\$0.0078$} & \shortstack{$-31.5$\\$0.078$} \\
Thompson bandit & GA & \shortstack{$-16.2$\\$0.742$} & \shortstack{$+5.9$\\$0.945$} & \shortstack{$-18.0$\\$0.028$} \\
 & RS & \shortstack{$+0.9$\\$0.844$} & \shortstack{$+3.4$\\$1.000$} & \shortstack{$-35.2$\\$0.039$} \\
\bottomrule
\end{tabular}
\end{table}

\textbf{Non-LLM preliminary conditions.} The archive-guided \texttt{RuleBasedProposer} is a stand-in, not evidence about an LLM. At $B=400$ on 1,200 candidates, random search nearly exhausts the frontier (chemotype-cold hypervolume standard deviation zero); all split-wise $p\ge0.17$. Reducing $B$ to 80 produces no advantage after correcting two bugs: median stand-in minus GA differences are $5.5/-11.9/0.4$, with $p=0.95/0.20/0.74$. Expanding to 30,044 candidates also yields no significant gain: $p=1.0/0.84/0.87$ versus GA and $0.74/0.46/0.74$ versus random search. These conditions establish that the negative findings are not solely saturation of the original library.

\begin{figure}[t]
\centering
\includegraphics[width=\linewidth]{Figure 2 LEAF.pdf}
\caption{Oracle validity, agentic value-add, and proposal diversity. \textbf{(A)} Held-out potency correlations for all three splits; the dashed line is $\rho=0.3$, and RAW264.7 potency on chemotype-cold is uninterpretable. \textbf{(B)} Random-split median paired hypervolume differences and 95\% bootstrap CIs for four models ($n=8$ each); asterisks denote Wilcoxon $p<0.05$. The sole positive point is Qwen2.5 14B versus static-GA ($+1.8$), not versus random search ($-4.3$); neither is significant. \textbf{(C,D)} Exact archive-duplicate proposal rates and minimum fragment-slot edit distances (0-3). Full-run aggregate duplicates are 71.5/67.0/3.5/32.3\% for Llama 8B/Qwen 7B/Qwen 14B/Gemini, versus 22.5/0.0\% for GA/random search; edit distances are 0.55/0.63/1.65/1.01 versus 1.52/2.81. Random search's zero is measured, not missing. Initial four-seed diagnostics are separately retained in Table~\ref{tab:original-4}.}
\label{fig:performance}
\end{figure}

\textbf{Four replicated LLMs.} Llama 3.1 8B and Qwen2.5 7B underperform numerically in all 12 split-by-baseline comparisons; four Llama and three Qwen comparisons reach $p<0.05$ (not five and three, as an earlier summary incorrectly stated). Llama's median deficits are $47.1/24.8/31.2$ versus GA and $42.0/28.8/46.3$ versus random search. Qwen 7B's corresponding deficits are $40.9/21.6/29.2$ and $29.1/34.3/43.4$. Complete test values are preserved in Tables~\ref{tab:original-1}-\ref{tab:original-2}; Figures~\ref{fig:performance}B and \ref{fig:diagnostics}D-E show the broader pattern.

Qwen2.5 14B, run on an RTX A4000 (16\,GB VRAM), is indistinguishable from both baselines on random and scaffold-cold splits: median differences $+1.8/+18.3$ versus GA and $-4.3/+33.9$ versus random search. On chemotype-cold it is worse by $55.1/77.0$ ($p=0.0078$ against both), with hypervolume confined to 84.7-84.9 across eight seeds. This coincides with the weakest oracle-validity regime; nearly flat oracle responses are a plausible explanation, not a demonstrated cause. Gemini Flash-Lite underperforms both baselines on all three splits ($p=0.0078$ throughout): differences $-91.8/-66.6/-24.8$ versus GA and $-85.5/-71.1/-60.2$ versus random search, with 0\% rejection throughout (Tables~\ref{tab:original-3}, \ref{tab:original-8}).

The random-split pooled analysis combines 32 (model, seed) differences: versus GA, 24/32 are negative (median $-38.7$); versus random search, 27/32 are negative (median $-32.6$). We report only these descriptive counts and medians here, not a $p$-value or confidence interval on the pooled sample: shared seeds and shared baseline runs across models violate the independence assumption that both bootstrap resampling and Wilcoxon signed-rank require, so any inferential statistic computed on this pooled set (a previous draft of this paper reported $p=8.2\times10^{-5}$ and $p=2.3\times10^{-6}$ here) would overstate precision rather than correct for the dependence, and we withdraw those numbers rather than retain a formally unjustified statistic alongside a caveat about it. This is descriptive corroboration that the direction is consistent across models, not a random-effects meta-analysis or any other dependence-aware inferential claim.

\textbf{Additional proposal controls.} The Thompson bandit is indistinguishable from both baselines on random/scaffold-cold splits, but worse on chemotype-cold ($-18.0$, $p=0.028$ versus GA; $-35.2$, $p=0.039$ versus random search; Table~\ref{tab:original-7}). This non-LLM result cautions against attributing cold-split difficulty exclusively to agentic reasoning. Llama 3.3 70B via Groq supplies only three random-split seeds: hypervolumes 107.1/116.2/116.2 versus GA's 182.1/112.5/135.5 and random search's 139.1/155.1/166.7. Median differences are $-19.3$ (CI $[-75.0,3.7]$) and $-38.8$ (CI $[-50.5,-32.0]$); these are exploratory, not significance claims. Rejection is zero; duplication is only 7.1\% and mean edit distance 1.76. Quota-constrained prompting exposes a sampled vocabulary subset per call, an additional protocol difference documented in Appendix~\ref{app:hosted}.

\subsection{Duplication, interventions, and practical cost}
Initial one-split, four-seed diagnostics show 67.8\%/62.5\% archive duplicates for Llama 8B/Qwen 7B and mean edit distances 0.587/0.650, compared with 22.5\%/1.516 for GA and 0\%/2.812 for random search. Llama labels 99.7\% of proposals ``crossover,'' Qwen 56.2\%; the baselines use different labels, so their reported 0\% is not absence of recombination. One real eight-candidate Llama pool repeats the identical triple six times (Appendix~\ref{app:example}). These are \emph{proposal} (pool-level) rates; external review correctly noted they must not be equated with selected-query fractions without directly measuring the latter. We instrumented the loop to log exactly this (Llama 8B, random split, $n=8$, $B=40$, $m=10$, $q=5$, 320 real charged queries): \textbf{79.1\%} of Llama 8B's \emph{acquisition-selected, budget-charging} queries exactly duplicate an already-evaluated archive member -- higher than the 67.8\% pool-level rate, not lower, meaning UCB acquisition does not filter duplicates out of the pool but if anything concentrates on them (plausible mechanism: a duplicate of a known-good archive member inherits that member's already-high posterior mean under $\mathcal O_{\mathrm{AL}}$, making it attractive to a value-seeking acquisition function with no explicit novelty term). The matched baselines' selected-query duplicate rates are 15.3\% (static-GA) and 0.0\% (random search, exact), both close to their own pool-level rates. This directly confirms budget waste as measured fact rather than an inferred opportunity: nearly four in five of Llama 8B's real oracle queries in this condition purchase no new information. Disabling UCB and querying the first $q$ surviving candidates produces no measurable improvement, consistent with poor pool diversity rather than selection alone; the selected-query result above clarifies that acquisition is, if anything, a second contributor to the deficit rather than a mitigating factor.

\textbf{Applicability domain of queried candidates.} The same instrumented run also tags every charged query by how many of its three fragment slots (A, B, C) hold a value absent from that split's real oracle \emph{training} partition -- 0 means a fully interpolated query within the labeled factorial design, 3 means the oracle must extrapolate on every slot. Across Llama 8B's 320 real queries, 10.3\% are fully in-domain (0 unseen values), 88.75\% have exactly one unseen value, 0.9\% have two, and none have all three unseen -- the oracle is never asked to extrapolate on every axis simultaneously for this model. Static-GA and random search extrapolate \emph{more}, not less, on average (mean unseen-value count 1.58 and 1.81 respectively, vs.\ Llama 8B's 0.91): both draw uniformly from the full 30,044-candidate vocabulary (dominated by the 62 added aldehyde values, \S3.3), while the LLM Proposer's archive-anchored crossover/mutation keeps it disproportionately close to the 1,200 originally labeled combinations. This means the LLM Proposer's measured deficit is not attributable to it querying more out-of-domain, harder-to-trust oracle predictions than the baselines -- if anything, the opposite is true, and its lower hypervolume must be explained by proposal duplication and lack of novelty (above) rather than the oracle penalizing its extrapolation risk. Coverage here is one model, one split, one 8-seed sweep; the same accounting for the other three fully-powered models and all three splits was not run given real per-seed compute cost (\S5) and is left as a direct, mechanically identical extension.

Qwen 14B does not share this failure: duplicate rates are 4.2\% on random/scaffold-cold and 2.1\% on chemotype-cold, with edit distances 1.63 and 1.69. Gemini's split-wise duplicates are 43.1/34.5/19.2\%, with distances 0.90/0.98/1.14. Aggregate cross-model differences do not imply a seed-level dose response: across 60 diagnosed runs, duplicate rate versus deficit has $\rho=-0.117$, $p=0.37$ (Figure~\ref{fig:diagnostics}C). Which novel molecules are proposed, acquisition, and oracle extrapolation remain possible sources of variation.

\begin{figure}[tbp]
\centering
\includegraphics[width=\linewidth]{Figure 3.pdf}
\caption{Learning curves, cost, and complete comparisons. \textbf{(A)} Hypervolume mean $\pm$ standard deviation: both baselines overtake Llama 8B within 2-3 rounds (random split, $n=8$). \textbf{(B)} Log-scale 80-call workload times (Table~\ref{tab:original-6}), approximately 300-800 times the non-LLM cost. \textbf{(C)} Negative seed-level robustness check ($n=60$, $\rho=-0.12$, $p=0.37$). \textbf{(D,E)} All 36 method/split/baseline comparisons ($n=8$ each); asterisks mark unadjusted Wilcoxon $p<0.05$. Positive cells are non-significant. The bandit is non-agentic; Groq's three-seed results are reported separately.}
\label{fig:diagnostics}
\end{figure}

A diversity-aware prompt explicitly lists triples to avoid. At $n=8$ on all splits, duplication falls from 67.8\% to 34.0\% and edit distance rises from 0.59 to 1.06. All six median comparisons remain negative: $-32.5/-15.6/-10.1$ versus GA and $-25.4/-18.0/-31.5$ versus random search. Only random-versus-random-search ($p=0.016$) and scaffold-versus-random-search ($p=0.0078$) remain significant.

\textbf{Primary vs.\ exploratory family (external review).} We declare the eight comparisons in Figure~\ref{fig:value-add} (the four fully-powered $n=8$ models $\times$ two baselines, random split) as this paper's single primary confirmatory family, subject to the Holm correction in Appendix~\ref{app:equivalence}. Every other reported comparison in this paper -- the diversity prompt, the Thompson bandit, all scaffold-cold/chemotype-cold splits, the temperature sweep, Groq's three seeds, the EHVI check, and both de-duplication checks -- is secondary or exploratory: individually informative and reported with exact statistics wherever computable, but not part of the family a reader should treat as confirmatory evidence for or against H1. The paper's overall conclusion is accordingly framed as \emph{no demonstrated improvement under the tested protocol}, not as multiplicity-corrected proof of individual-model inferiority beyond the primary family, and not as evidence of equivalence anywhere (Appendix~\ref{app:equivalence}).

A four-seed temperature check is exploratory: temperature 0.2 produces 87.5\% duplicates and identical hypervolume 106.4 on every seed; temperature 1.0 yields 72.2\% duplicates and median hypervolume 126.3 versus 110.4 at matched default-temperature seeds. Its duplicate rate remains above the separately reported default diagnostic rate; the comparison does not establish monotonic temperature improvement.

The matched \texttt{dedupe\_pool} intervention removes exact archive matches \emph{before} surrogate scoring for every method, without backfill. An initial reduced-scope check (random-split Llama 8B, $B=15$, $n=4$) shifted the median difference versus GA from $-17.6$ to $-4.7$ and versus random search from $-15.1$ to $+2.3$, with all CIs crossing zero (Table~\ref{tab:original-5}). This was re-run at full statistical power ($B=40$, $n=8$, random split, matching the primary protocol exactly), a direct response to external review requesting the matched check at $n\ge8$ rather than $n=4$, and subsequently extended to the other two splits at $n=4$ (exploratory scope: the random-split $n=8$ dedup arm alone measured $\approx$10 hours of real wall-clock time across its eight seeds, making an $n=8$, two-more-split extension a multi-day commitment not undertaken here):

\begin{center}
\begin{tabular}{@{}lcccc@{}}
\toprule
Split & $n$ & No dedup (control) & With dedup \\
\midrule
Random & 8 & $-33.0$ / $-38.7$ ($p=0.023$/$0.0078$) & $-21.9$ / $-35.7$ ($p=0.195$/$0.055$) \\
Scaffold-cold & 4 & $-34.2$ / $-40.0$ (unreliable) & $-5.1$ / $-13.4$ (unreliable) \\
Chemotype-cold & 4 & $-45.3$ / $-75.1$ (unreliable) & $-35.0$ / $-70.2$ (unreliable) \\
\bottomrule
\end{tabular}
\end{center}
{\small Each cell: median paired difference vs.\ static-GA / vs.\ random search.}

At full power (random split), de-duplication narrows both medians and moves both comparisons from significant to non-significant (the random-search comparison only marginally, $p=0.055$); it does not reverse either median's sign, unlike the smaller $n=4$ pilot's random-search flip, which does not survive the larger sample. The control condition itself is an independent full-power replication of the primary random-split result (original: $-47.1$/$-42.0$; here: $-33.0$/$-38.7$), same direction and significance pattern but different absolute magnitude -- attributable to genuine run-to-run variance in the real LLM Proposer's stochastic output, not a discrepancy requiring correction. The two $n=4$ extension splits show a real, striking asymmetry rather than a uniform effect: scaffold-cold's deficit shrinks dramatically under de-duplication (roughly 6-fold vs.\ GA, 3-fold vs.\ random search), closely mirroring the random split's direction, while chemotype-cold's deficit barely moves (shrinking by only about a fifth to a quarter). This is consistent with chemotype-cold's independently-flagged status (\S6.1) as the split with the weakest oracle validity and a different failure signature at 14B (\S6.2): duplicate-driven waste appears to be a smaller share of the deficit specifically on this split, suggesting a mechanism beyond proposal duplication dominates there, not resolved by this check alone. Spending the same budget under dedup took substantially longer than the pilot's 2.3--7$\times$ multiplier on the random split: two of eight seeds (4 and 7) never spent the full $B=40$ budget at all (31/40 and 14/40 queries respectively) before triggering the loop's excessive-empty-rounds safety valve, because real proposal duplication was severe enough in those seeds that no novel candidate pool could be assembled; this did not recur in the $n=4$ scaffold-cold/chemotype-cold extension (all seeds spent the full budget). This is itself a finding, not an artifact to explain away: de-duplication does not only trade oracle queries for proposal rounds, it can make the nominal budget partially or fully unspendable for a real LLM Proposer, at least on some splits. This controlled check supports duplicate-driven waste as a real, substantial, but split-dependent component of the deficit, without establishing that the LLM Proposer outperforms either baseline once duplicates are removed.

\textbf{Stricter de-duplication: excluding all previously evaluated candidates, not only the archive (external review item \#3).} The check above excludes exact matches of the current \emph{Pareto-optimal} archive only; a dominated-but-already-evaluated candidate could still be re-proposed and re-queried, and within-pool duplicates (two identical novel proposals in the same round) were not addressed. We implemented a stricter, more literal control (\texttt{dedupe\_scope=all\_evaluated}): excludes exact matches against \emph{every} candidate ever queried this run, dominated or not, plus within-pool duplicates, with no backfill in either scope (consistent with this paper's existing no-backfill design; the outer round loop is the replenishment mechanism in both cases). At $n=4$ (random split, matching seeds 0--3 of the existing control for a direct comparison):

\begin{center}
\begin{tabular}{lcc}
\toprule
Comparison & No dedup (control, same 4 seeds) & Stricter dedup (all evaluated) \\
\midrule
vs.\ static-GA & median $-33.5$ & median $\mathbf{+20.1}$ \\
vs.\ random search & median $-52.5$ & median $-15.4$ \\
\bottomrule
\end{tabular}
\end{center}

Under this stricter criterion, the median difference against static-GA \emph{flips sign} and the deficit against random search shrinks by roughly 70\%. Both new-condition CIs are wide and cross zero ($[-46.6, 119.8]$ vs.\ GA; $[-32.2, 77.2]$ vs.\ RS), and $n=4$ cannot support a reliable significance claim (\S4.3) -- this is not evidence that the LLM Proposer outperforms static-GA under this protocol, only that duplicate-driven waste, measured at its most literal and complete definition, accounts for enough of the deficit that removing it can plausibly reverse the median at this small sample. This is a stronger, more favorable-to-the-Proposer result than the archive-only check above, consistent with archive-only de-duplication being an incomplete novelty control, as external review noted. A full $n\ge8$, all-split version of this stricter check is the natural next step and is not run here given the real per-seed compute cost already established for the archive-only version.

\textbf{Richer prompt with real objective scores (external review item \#2).} The default prompt (Appendix~\ref{app:llm-repro}) shows only fragment identities, never the actual potency values or an objective description -- a real, previously-unaddressed limitation on what the LLM Proposer is asked to optimize toward beyond archive membership. We implemented a \texttt{v3\_scored} prompt variant that additionally shows each sampled archive member's real \texttt{potency}/\texttt{potency\_raw} scores and an explicit maximize-both-objectives description. At $n=4$ (random split, otherwise identical protocol):

\begin{center}
\begin{tabular}{lcc}
\toprule
Comparison & Default prompt (same 4 seeds) & With real objective scores \\
\midrule
vs.\ static-GA & median $-33.5$ & median $\mathbf{+13.9}$ \\
vs.\ random search & median $-52.5$ & median $-32.7$ \\
\bottomrule
\end{tabular}
\end{center}

Adding real objective scores also flips the median against static-GA positive and narrows the random-search deficit by roughly 38\%, at $n=4$ with wide, zero-crossing CIs ($[-63.2,119.8]$ vs.\ GA; $[-42.9,77.2]$ vs.\ RS) -- directional, not a significance claim. \textbf{A caveat we verified directly rather than let pass silently:} seed 1's hypervolume under this richer prompt (232.263) is numerically identical, to twelve significant figures, to seed 1's hypervolume under the independent stricter-dedup check above. Inspecting both runs' real per-query fragment logs (\texttt{queries.csv}) confirms the two runs are not a caching artifact -- their query trajectories genuinely differ round by round -- but both independently converge on an equally strong final candidate from the same seed-determined starting archive, giving this one seed disproportionate leverage on both reported medians. This is consistent with, not contradictory to, this paper's broader finding (\S3.2) that the real oracle's non-dominated frontier concentrates on a small number of attainable values; it means both of these $n=4$ exploratory results should be read as weaker evidence than four independent draws would suggest, not stronger. Combining a richer prompt with de-duplication, and repeating both at full statistical power, are natural next steps beyond this paper's scope.

\textbf{Leakage probe (external review): does the richer-prompt gain reflect memorized chemistry knowledge rather than genuine use of the injected scores?} Showing the LLM real potency values in \texttt{v3\_scored} raises a real confound: an improvement could come from the model reactivating pretrained associations between similar structures and potency, not from reasoning over the specific scores given in that call. We ran the same bare-LLM leakage probe design used for H4 (\S2.3), applied to this confound directly: \texttt{llama3.1:8b} -- the identical model and backend used as the Proposer throughout -- was asked to zero-shot-score 50 real, held-out random-split AGILE lipids' transfection efficiency (1--10 scale, temperature 0, bare SMILES string only, no archive, no vocabulary, no scores), the same protocol already used for the Qwen3-8B/LipoAgent comparison (\S3.1). The result is more degenerate than a low correlation: \textbf{the model returned the identical score (8.0) for all 50 structurally distinct molecules}, a constant, non-discriminating output for which a correlation coefficient is undefined rather than merely small. This is a stronger result against the memorization explanation than a low-but-nonzero correlation would have been: the bare model exhibits no exploitable zero-shot signal whatsoever about real potency, degenerate rather than weakly informative, making it implausible that \texttt{v3\_scored}'s improvement reflects reactivated pretrained knowledge about which specific structures are potent. The more plausible explanation is that the model is doing something with the real scores it is shown in-context, though this probe cannot distinguish genuine structure-activity reasoning from simpler explanations (e.g.\ anchoring on and interpolating between the numeric values shown) -- disentangling those is future work (rank-only context, below).

\textbf{Rank-only (ordinal) context: is magnitude necessary, or does relative preference alone explain the effect?} We implemented a fourth prompt variant, \texttt{v4\_rank}, that shows each sampled archive member's rank (1 = best of the up-to-8 shown, by summed real potency scores) instead of its raw values -- verified directly on a synthetic example that no magnitude or scale information leaks into the rendered prompt. Run under the identical exploratory protocol ($n=4$, random split, same four seeds as \texttt{v3\_scored} for a direct comparison):

\begin{center}
\begin{tabular}{lccc}
\toprule
Comparison & Default prompt & \texttt{v3\_scored} (raw values) & \texttt{v4\_rank} (rank only) \\
\midrule
vs.\ static-GA & median $-33.5$ & median $+13.9$ & median $\mathbf{+17.2}$ \\
vs.\ random search & median $-52.5$ & median $-32.7$ & median $\mathbf{-17.7}$ \\
\bottomrule
\end{tabular}
\end{center}

Rank-only context matches or exceeds the raw-values result on both comparisons at this sample size (CIs remain wide and zero-crossing for both enriched-context variants; $[-31.0,114.8]$ vs.\ GA and $[-34.2,72.2]$ vs.\ RS for \texttt{v4\_rank}) -- relative preference alone accounts for at least as much of the apparent improvement as exact magnitudes do, which weighs against a magnitude-anchoring/interpolation explanation and is more consistent with the model using the injected signal as a genuine preference ordering. Seed 1 again produces the run's strongest hypervolume (227.271), close to but not identical to \texttt{v3\_scored} and the stricter-dedup check's seed-1 values (232.263 in both) -- a third independent condition converging near the same high-value region from the same seed-determined starting archive, further supporting the earlier explanation (\S3.3) that this reflects a real, discoverable attractor in the search space rather than a run-to-run artifact, and reinforcing that these $n=4$ results should be read cautiously given the shared seed's outsized influence. Both enriched-context variants remain exploratory, directional findings, not significance claims; a full $n\ge8$ sweep of both is the natural next step.

\textbf{Score-shuffle perturbation (\texttt{v6\_shuffled}): does the result depend on the true score-to-fragment mapping?} Following a diagnostic design used to test whether biological-reasoning LLMs actually use the inputs they are given (shuffling an input and checking whether the output changes) \citep{fang2026bioreasoninputs}, we built the direct analog for \texttt{v3\_scored}: an otherwise identical prompt in which the same real score values are randomly reassigned to the wrong fragments, breaking the true correspondence while keeping every other surface feature (phrasing, objective description, numeric format) identical. At $n=4$, same seeds: median $+13.9$ vs.\ static-GA and $-32.7$ vs.\ random search -- statistically indistinguishable from \texttt{v3\_scored}'s own $+13.9$/$-32.7$ at this sample size. Taken at face value this would replicate the reference study's finding that shuffling an input leaves model behavior unchanged. We checked the real per-query trajectories (\texttt{queries.csv}) before accepting that reading, and the mechanism is more specific: one of four seeds (seed 2) has an identical real fragment queried in all 40 of 40 rounds regardless of whether scores are real or shuffled -- a genuine, round-by-round null effect for that seed -- while the other three seeds' real per-round choices do diverge between conditions (11/40, 25/40, and 31/40 identical queries for seeds 0, 1, and 3 respectively), yet still converge on the same or a very similar final hypervolume in two of those three cases. This is consistent with, not contradictory to, this paper's own repeated finding (\S3.2, \S3.3) that the real attainable Pareto frontier here is narrow enough that different search trajectories often land on the same few endpoints -- a real, substantive limitation on what a final-hypervolume comparison at $n=4$ can distinguish, distinct from (and more specific than) a claim that the model ignores the score-to-fragment mapping outright. A trajectory-level analysis (e.g.\ comparing per-round proposal identity rather than only final hypervolume, at full statistical power) is needed to resolve how much of \texttt{v3\_scored}'s effect reflects genuine use of the mapping versus this frontier-concentration effect, and is left for future work.

\textbf{Evidence-conflict probe: does the model follow a (possibly false) injected quality signal?} The reference study's second diagnostic tool makes an injected input disagree with another available signal and measures which one the model follows \citep{fang2026bioreasoninputs}. We built the direct analog as a standalone, controlled probe (not embedded in the search loop, for a cleaner signal): a fixed, real 8-member sample from the labeled AGILE set, rank-labeled exactly as in \texttt{v4\_rank} except the ranks are \emph{deliberately inverted} -- the member shown as ``rank 1 (best)'' is the real worst of the eight by true summed score, and vice versa. Llama 3.1 8B (the same model used throughout) was queried 40 times against this fixed, falsely-labeled archive; each returned candidate's nearest neighbor among the eight shown members was identified by fragment-slot edit distance (13 of 40 calls were ties or unparseable and excluded). Of the 27 resolved calls, only 14.8\% land nearest the falsely-labeled ``best'' member -- statistically indistinguishable from the 12.5\% rate expected if rank were ignored entirely -- and \emph{zero} land nearest the real best member. Instead, 74\% (20 of 27) land nearest whichever fragment happened to be listed \emph{first} in the archive, regardless of its real or falsely-labeled rank. This is a sharper finding than either "follows the false label" or "ignores the label": in this probe, Llama 3.1 8B's proposals are dominated by list position, not by the quality signal in either its true or falsified form. We disclose a real limitation of this probe's design rather than overstate the result: the eight-member archive and its listing order were held fixed across all 40 calls to isolate the rank-label manipulation, which means list position and fragment identity are perfectly confounded here -- a follow-up that additionally randomizes listing order per call is needed to confirm position dominance over content-driven preference for this specific fragment, and is left for future work. The finding is, however, directly consistent with this paper's independently-established duplicate-heavy, low-diversity proposal mechanism (\S3.3): whatever drives Llama 3.1 8B's choice of which archive member to build from, in this controlled probe it is evidently not the explicit quality ranking it is shown, true or false.

\textbf{Genuine reasoning-mode Proposer.} The checks above test an LLM prompted for a rationale (\texttt{v5\_cot}) but never a model with native extended inference-time reasoning -- a real, distinct mechanism (Ollama's \texttt{think} API for \texttt{qwen3:8b}; verified directly that forcing schema-constrained JSON output suppresses the reasoning phase entirely, so this condition runs without \texttt{format="json"}, relying on the model's own reliably well-formed output instead, with a graceful fallback for the rare malformed response). Under the default prompt (no scores shown), $n=4$, random split: median $-26.6$ vs.\ static-GA (CI $[-57.4,24.1]$, crossing zero) and $-55.0$ vs.\ random search (CI $[-73.4,-18.4]$, not crossing zero) -- directionally similar to the non-reasoning models (\S3.2), not a reversal, and if anything numerically worse against random search than Llama 3.1 8B's own default-prompt result ($-42.0$). Genuine reasoning traces were captured for 155 of 160 real queries (mean length 5{,}669 characters -- substantial, not token-limited), but \textbf{zero of 155 traces contain the word ``potency''} despite the model producing thousands of characters of reasoning per call. Under this prompt the objective is never named or described (\S2.2, Appendix~E), so this is consistent with the model reasoning about structural plausibility rather than the (unstated) optimization target, not yet evidence that a reasoning-capable model ignores an objective it is actually given; the scored-prompt condition (identical reasoning setup, real scores and an explicit maximize-both-objectives instruction) is the condition that can distinguish these, and is reported next. Reasoning mode measured $\approx$61 minutes for one 40-query seed-run, roughly 5--10$\times$ slower than every non-reasoning model in Table~\ref{tab:original-6} -- a real, substantial practical cost for whatever benefit genuine reasoning provides here.

Under the real objective-scored prompt (\texttt{v3\_scored}, same reasoning setup, $n=4$, same seeds): median $-19.1$ vs.\ static-GA (CI $[-58.2,10.7]$, crossing zero) and $-53.4$ vs.\ random search (CI $[-56.4,-37.9]$, not crossing zero). The contrast with the default-prompt condition directly answers the question posed above: of 152 of 160 captured traces (mean length 6{,}976 characters, longer than the default condition's), \textbf{81.6\% explicitly mention ``potency''} -- a sharp, real increase from 0\% when the objective is actually named and real scores are provided, confirming the model genuinely engages with the injected evidence when it is given something concrete to engage with, unlike the uniform silence on the topic under the uninformative default prompt. This is a real point in favor of genuine use of the signal, independent of and complementary to the evidence-conflict probe above. It does not, however, translate into better search outcomes than the much cheaper non-reasoning condition: Llama 3.1 8B under the identical scored prompt reached a \emph{positive} median against static-GA ($+13.9$) and a smaller deficit against random search ($-32.7$, \S3.3), both better than the reasoning model's $-19.1$/$-53.4$ despite the reasoning model producing roughly 7{,}000 characters of substantive, evidently on-topic reasoning per call at 5--10$\times$ the latency. Genuinely engaging with the injected evidence, measured directly rather than inferred, is not sufficient to outperform a smaller non-reasoning model given the same evidence in this budget-matched setting -- a further, model-architecture-level instance of this paper's central finding that more apparent sophistication in the proposal mechanism does not reliably translate into matched-budget search advantage.

Per-candidate latencies are approximately 4.0/10.3/10.3/5.1 seconds for Llama 8B/Qwen 7B/Qwen 14B/Gemini. Groq inference takes about 0.6 seconds raw but is paced at about seven seconds; non-LLM search requires no LLM calls (Table~\ref{tab:original-6}). The original cost table reports under one second per non-LLM run and sub-second GP rounds. These workload-specific measurements should not be interpreted as universal end-to-end timings; hosted quotas, retries, and multi-day accumulation add further cost beyond nominal 80-call times.

\subsection{Secondary-agent ablations and implementation audit}
\textbf{Verifier.} The Verifier was enabled in every experiment reported in this paper (\texttt{use\_verifier: true}). Because AGILE contains no experimental toxicity labels, the conformal calibration set satisfies \texttt{cal\_scores = cal\_labels} exactly, producing zero-width intervals. Consequently the Verifier reduces to a deterministic toxicophore filter and contributes no additional epistemic signal; binary-filter coverage is correspondingly undefined, not zero. The Verifier ablation is therefore uninformative about whether debate-style verification helps: the component was already degenerate before the ablation was run, not shown to be ineffective by it.

\textbf{Hypothesis Ledger.} The Ledger was enabled in the primary sweep and early ablations. Across all 72 primary runs (eight seeds, three splits, three methods) it remained empty: hypothesis precision is undefined, not zero. A reproduced $B=40$, $q=5$, $k=4$ run ends with 12 Pareto members sharing one isocyanide value, different from the fixed control used for pairing -- restricting matches to the Pareto archive prevents completed tests even though the validation logic passes unit tests in isolation. Under the shrunk-budget extended-vocabulary check ($B=80$, $m=40$, $q=8$, $n=8$), Ledger-disabled chemotype-cold results are byte-identical to full LEAF, consistent with an empty Ledger having nothing to feed back regardless. In every subsequent isolation experiment in this paper (de-duplication in all its variants, EHVI, the single-assay chemotype-cold reanalysis, and the three enriched-prompt checks), the Ledger was explicitly disabled (\texttt{use\_hypothesis: false}) rather than left to redundantly reconfirm the same empty result, to focus evaluation on the proposer and acquisition function. The Ledger ablation is uninformative about active hypothesis feedback for two distinct reasons depending on which runs it refers to: it never populated under the primary-sweep matching constraints, and it was not exercised at all in the later runs. Neither finding establishes that literature-grounded hypothesis feedback would fail to help under matching constraints that actually produce testable pairs.

An early stand-in result appeared strongly positive: median GA gap 92.5, CI $[30.1,119.1]$, $p=0.0078$. Two errors explained it: reusing a full-data oracle made split labels ineffective, and GA mutation probability 0.2 was unmatched to the stand-in's 50\% pure mutation. Fixing both erased the effect. A later audit found a third error: the acquisition configuration was read without dispatching its EHVI option. All existing runs actually used UCB, with configured $\beta=2.0$ matching the hardcoded value; the reported UCB comparisons survive. After fixing the dispatch, a first real, exploratory exercise of the EHVI path ($n=4$, random split, otherwise identical to the primary protocol) shows median differences of $-23.5$ versus static-GA and $-38.2$ versus random search; $n=4$ makes Wilcoxon unreliable (\S4.3), so this is directional, not a significance claim, but the magnitude and direction match the UCB result closely rather than suggesting EHVI reverses the pattern. A full $n\ge8$, all-split EHVI sweep remains future work. These corrections, and the de-duplication and EHVI checks, are part of the evaluation record rather than evidence for an improved lipid-design system.

\section{Discussion, limitations, and conclusion}
\label{sec:discussion}
No tested configuration establishes an agentic advantage under matched oracle-query budgets. Small-model duplication is a measured failure mode, but not a universal explanation: Qwen 14B and Groq produce diverse proposals without demonstrating superior search. Prompting mitigates repetition, and exploratory structural de-duplication reduces the deficit at substantial proposal-time cost. The relevant deployment trade-off is therefore query efficiency \emph{and} compute/latency, not hypervolume alone.

The oracle is a fitted model trained on public assays, not a laboratory. Its bias can yield self-consistent but biologically false frontiers or hypothesis ledgers; weaker chemotype-cold validity complicates interpretation even above the gate. The conditional synthesizability objective may be degenerate, toxicity is a rule rather than an assay, and tropism heuristics would circularly rediscover their construction. Exact objective/reference-point provenance and access to the reported configurations are necessary for reproducibility. The 50-molecule scoring comparison is limited to one Qwen3 component, not LipoAgent's full system. Reported bootstrap and unadjusted Wilcoxon results are small-sample summaries, not equivalence tests or multiplicity-corrected population-wide evidence; the pooled analysis does not model shared baselines or model heterogeneity. We checked both concerns directly rather than leaving them as caveats (Appendix~\ref{app:equivalence}): two one-sided equivalence tests at a $\pm10\%$-of-baseline-median threshold find \emph{no} primary comparison, including Qwen2.5 14B's nominally ``indistinguishable'' random- and scaffold-cold results, reaches equivalence significance at $n=8$ -- the honest statement is no evidence of a difference and no evidence of equivalence, not demonstrated parity. Holm-Bonferroni correction across the eight-comparison primary random-split family leaves only one comparison significant in the strict sense, though three sit tied at Wilcoxon's $n=8$ minimum achievable $p$-value and which one survives correction is an artifact of tie-breaking, not a meaningful distinction between them.

Two secondary components were architecturally present but empirically non-contributory. The Verifier lacked a real toxicity signal (AGILE contains no experimental toxicity labels, \S3.4), and the Hypothesis Ledger never populated under the primary sweep's Pareto-restricted matching constraints. Later isolation experiments disabled the Ledger deliberately to focus evaluation on the proposer. These facts limit any claim about the value of the full multi-agent design: they show what did not happen under the tested conditions, not that the components cannot contribute under different ones (e.g.\ a real toxicity assay, or counterfactual matching that is not Pareto-restricted).

Coverage spans four eight-seed models on three splits and one replicated prompt variant, not arbitrary architectures or prompts. The 70B, temperature, and de-duplication checks are underpowered; prompt compression further limits direct attribution of Groq's outcome to model scale. Meditron 7B fails the required JSON/vocabulary schema (100\% rejection in repeated checks), so no misleading fallback-dominated sweep is reported. TxGemma access and 27B memory requirements, unavailable prebuilt chemistry models, and hosted quotas constrain broader evaluation. The recorded Groq catalog change is an additional historical availability constraint, not a claim that the model remains unavailable today.

Priority extensions are matched de-duplication and EHVI sweeps at $n\ge8$ across all splits; independently labeled toxicity and organ-selectivity data; counterfactual controls drawn beyond the concentrated Pareto set; a full prompt/temperature study; and trained generative or conditional-MOBO baselines. Halogenated, literature-sourced Ugi-compatible reagents are needed before H4 can run. External LUMI-lab validation, 3D/membrane-binding information, cross-barrier transfer, and out-of-distribution/adversarial tests remain future work, not results. Richer oracles and VAE/diffusion proposers require separate validation rather than assuming that scaling the same agentic loop will help.

\textbf{Ethics and broader impact.} This computational study uses public data and neither synthesizes nor experimentally administers compounds. Its vocabulary is restricted to literature-precedented ionizable-lipid routes, without unrestricted de novo scaffold generation. Rule-based screening does not establish safety, and none of the reported candidates should be construed as clinically validated.

\textbf{Conclusion.} LEAF contributes a controlled test of whether an agentic component earns its cost. The honest summary is \emph{no demonstrated improvement under the tested protocol}: no primary comparison establishes an agentic advantage, none establishes equivalence to baseline either, and several models incur a reproducible deficit, partly explained by duplicate proposals. Transparent nulls, exploratory interventions, and corrected false positives are the result, not a claim that LEAF designs better lipids than simpler methods, and not a claim of multiplicity-corrected evidence of individual-model inferiority beyond the primary family declared in \S6.2.

\label{main-text-end}
\clearpage
\ifnum\value{page}>10
  \PackageError{LEAF}{Main content exceeds nine pages}{Condense the main text or reposition floats without changing the NeurIPS style.}
\fi
\begingroup
\small
\begin{thebibliography}{99}
\label{references-start}
\bibitem[Angelopoulos and Bates(2023)]{angelopoulos2023}
A.~N. Angelopoulos and S.~Bates.
A gentle introduction to conformal prediction and distribution-free uncertainty quantification.
\emph{Foundations and Trends in Machine Learning}, 2023.

\bibitem[Boiko et al.(2023)]{boiko2023coscientist}
D.~A. Boiko, R.~MacKnight, B.~Kline, and G.~Gomes.
Autonomous chemical research with large language models.
\emph{Nature}, 624:570--578, 2023.

\bibitem[Bran et al.(2024)]{bran2023chemcrow}
A.~M. Bran, S.~Cox, O.~Schilter, C.~Baldassari, A.~D. White, and P.~Schwaller.
Augmenting large language models with chemistry tools.
\emph{Nature Machine Intelligence}, 6:525--535, 2024.

\bibitem[Fang et al.(2026)]{fang2026bioreasoninputs}
Z.~Fang, J.~Thoduguli, A.~Fesser, R.~Zhang, S.~Kakade, and M.~Zitnik.
When do biological reasoning models use their biological inputs?
\emph{arXiv preprint arXiv:2610.00898}, 2026.

\bibitem[Gessner et al.(2024)]{gessner2024}
A.~Gessner, S.~W. Ober, O.~Vickery, D.~Ogli\'c, and T.~U\c{c}ar.
Active learning for affinity prediction of antibodies. \emph{arXiv preprint arXiv:2406.07263}, 2024.

\bibitem[Gorantla et al.(2024)]{gorantla2024benchmarking}
R.~Gorantla, A.~Kubincov\'a, B.~Suutari, B.~P. Cossins, and A.~S.~J.~S. Mey.
Benchmarking active learning protocols for ligand-binding affinity prediction.
\emph{Journal of Chemical Information and Modeling}, 64(6):\allowbreak1955--1965, 2024.
doi:10.1021/acs.jcim.4c00220.

\bibitem[Gottweis et al.(2026)]{gottweis2026}
J.~Gottweis et al. Accelerating scientific discovery with Co-Scientist.
\emph{Nature}, 2026. doi:10.1038/s41586-026-10644-y.

\bibitem[Ghafarollahi and Buehler(2024)]{protagents2025}
A.~Ghafarollahi and M.~J. Buehler.
ProtAgents: Protein discovery via large language model multi-agent collaborations combining physics and machine learning.
\emph{arXiv preprint arXiv:2402.04268}, 2024.

\bibitem[Li et al.(2026)]{li2025lipoagent}
L.~Li, A.~Lu, H.~Wang, Z.~Feng, C.~Duan, Q.~Bao, Z.~Zhao, and S.~Q. Zhang.
LipoAgent: Coordinating fine-tuned LLM agents for safer lipid design.
\emph{Findings of the Association for Computational Linguistics: ACL 2026}, pages 40070--40081, 2026.
doi:10.18653/v1/2026.findings-acl.1992.

\bibitem[Lundberg and Lee(2017)]{lundberg2017}
S.~M. Lundberg and S.-I. Lee.
A unified approach to interpreting model predictions.
\emph{NeurIPS}, 30, 2017.

\bibitem[Xu et al.(2026)]{lumilab2026}
Y.~Xu, H.~Cui, K.~Pang, G.~Li, F.~Gong, S.~Dong, B.~Wang, and B.~Li.
LUMI-lab: A foundation model-driven autonomous platform enabling discovery of ionizable lipid designs for mRNA delivery.
\emph{Cell}, 189(6):1620--1635.e25, 2026. doi:10.1016/j.cell.2026.01.012.

\bibitem[Rasmussen and Williams(2006)]{rasmussen2006}
C.~E. Rasmussen and C.~K.~I. Williams.
\emph{Gaussian Processes for Machine Learning}. MIT Press, 2006.

\bibitem[Singh et al.(2026ppi)]{singh2026explainable}
H.~Singh, R.~K. Singh, S.~P. Srivastava, S.~Pradhan, and R.~Gorantla.
Explainable protein-protein binding affinity prediction via fine-tuning protein language models.
\emph{bioRxiv preprint}, 2026.
doi:10.64898/2026.03.30.715237.

\bibitem[Singh et al.(2026antibody)]{singh2026antibody}
H.~Singh, A.~Malhotra, S.~P. Srivastava, R.~K. Singh, and R.~Gorantla.
Antibody-antigen affinity prediction with chain-aware protein language modeling.
\emph{bioRxiv preprint}, 2026.
doi:10.64898/2026.06.19.733375.

\bibitem[Srinivas et al.(2009)]{srinivas2009}
N.~Srinivas, A.~Krause, S.~M. Kakade, and M.~Seeger.
Gaussian process optimization in the bandit setting: No regret and experimental design. 2009.

\bibitem[Srivastava et al.(2026)]{srivastava2026explainable}
S.~P. Srivastava, R.~Gorantla, S.~K. Chundru, C.~J.~R. Winkelman, A.~S.~J.~S. Mey, and R.~K. Singh.
Explainable active learning framework for ligand binding affinity prediction.
\emph{Digital Discovery}, 5(2):769--779, 2026.
doi:10.1039/d5dd00436e.

\bibitem[Settles(2012)]{settles2012}
B.~Settles. \emph{Active Learning}. Morgan \& Claypool, 2012.

\bibitem[Sundararajan et al.(2017)]{sundararajan2017}
M.~Sundararajan, A.~Taly, and Q.~Yan.
Axiomatic attribution for deep networks. \emph{ICML}, 2017.

\bibitem[Swanson et al.(2025)]{swanson2025virtuallab}
K.~Swanson, W.~Wu, N.~L. Bulaong, J.~E. Pak, and J.~Zou.
The Virtual Lab of AI agents designs new SARS-CoV-2 nanobodies.
\emph{Nature}, 646:716--723, 2025. doi:10.1038/s41586-025-09442-9.

\bibitem[Wang et al.(2025)]{swarmdesign2025}
F.~Y. Wang, D.~S. Lee, D.~L. Kaplan, and M.~J. Buehler.
Swarms of large language model agents for protein sequence design with experimental validation.
\emph{arXiv preprint arXiv:2511.22311}, 2025.

\bibitem[Zhou et al.(2026)]{toolmol2026}
A.~Y. Zhou, S.~Vadgama, S.~Varambally, P.~Eckmann, M.~K. Gilson, and R.~Yu.
ToolMol: Evolutionary agentic framework for multi-objective drug discovery.
\emph{arXiv preprint arXiv:2605.12784}, 2026.

\bibitem[Chen et al.(2023)]{xu2023ugi}
J.~Chen, Y.~Xu, M.~Zhou, S.~Xu, A.~J. Varley, A.~Golubovic, R.~X.~Z. Lu, K.~C. Wang, M.~Yeganeh, D.~Vosoughi, and B.~Li.
Combinatorial design of ionizable lipid nanoparticles for muscle-selective mRNA delivery with minimized off-target effects.
\emph{Proceedings of the National Academy of Sciences}, 2023. doi:10.1073/pnas.2309472120.

\bibitem[Xu et al.(2024)]{xu2024agile}
Y.~Xu, S.~Ma, H.~Cui, J.~Chen, S.~Xu, F.~Gong, A.~Golubovic, M.~Zhou, K.~C. Wang, A.~Varley, R.~X.~Z. Lu, B.~Wang, and B.~Li.
AGILE platform: A deep learning powered approach to accelerate LNP development for mRNA delivery.
\emph{Nature Communications}, 2024. doi:10.1038/s41467-024-50619-z.

\bibitem[Ou et al.(2024)]{zhu2024generative}
Y.~Ou, J.~Zhao, A.~Tripp, M.~Rasoulianboroujeni, and J.~M. Hern\'andez-Lobato.
A deep generative model for the design of synthesizable ionizable lipids.
\emph{NeurIPS Workshop on AI for New Drug Modalities}, 2024. arXiv:2412.00928.
\end{thebibliography}
\endgroup
\clearpage
\appendix
\section{Secondary-component specification and implementation scope}
\label{app:agents}
\subsection{Debate and conformal verification}
The general Verifier comprises $N$ independently prompted toxicologists spanning at least two base models. Each uses RDKit toxicophore matching, an active-learning toxicity posterior when available, and free-text reasoning. Same-model prompt or temperature variants measure prompt sensitivity rather than independent epistemic diversity. Inter-model agreement above 95\% flags insufficient diversity for H3.

Held-out split-conformal calibration converts aggregated scores to intervals \citep{angelopoulos2023}; intervals crossing the decision threshold trigger abstention. Adaptive acquisition and refitting invalidate ordinary exchangeability. Weighted or Mondrian recalibration therefore occurs each round, with empirical coverage monitoring rather than a distribution-free guarantee under arbitrary shift. Calibration provenance, weighting, strata, and thresholds must be specified.

The AGILE experiments contain no learned toxicity posterior. Calibration scores equal deterministic toxicophore labels, so residual uncertainty and interval width are zero. Repeated rule outputs reproduce binary-filter decisions, not independent epistemic evidence or demonstrated calibration of a learned safety signal. Testing this component requires independent toxicity labels and assay harmonization.

\subsection{Hypothesis generation and testing}
Every $k$ rounds, indexed LNP literature and Integrated Gradients or KernelSHAP attributions over the Pareto archive motivate structure--property hypotheses. Each rule requires matched counterfactual pairs, a minimum pair count, and a paired Wilcoxon test. Only frozen $\mathcal O_{\mathrm{GT}}$ supplies validation labels. All queried molecules, including controls, consume the shared budget. Incomplete pair sets remain untested; the Ledger updates only after validation and then informs proposals when enabled. Algorithm~\ref{alg:leaf} bounds complete pair sets by the remaining budget.

Pair construction searches only Pareto members and requires a fixed control reagent. In the reproduced $B=40$, $q=5$, $k=4$ configuration, all 12 final members share one isocyanide among 14 possible values, different from that control. Correlated objectives can concentrate the frontier and eliminate matched controls. This explains the empty Ledger and undefined precision across 72 original sweep runs. Passing validation-logic unit tests does not establish that the optimization loop supplies testable pairs.

The LUMI-lab association concerns brominated tails and potency, not organ tropism. Because its preprint was public from February 2025, a bare-LLM probe without retrieval or archive context must precede any rediscovery claim, with either outcome disclosed. Detected prior knowledge changes the interpretation to consistency checking. Neither the original AGILE vocabulary nor its virtual-library extension supplies halogenated reagents. H4 therefore requires literature-sourced, Ugi-compatible halogenated reagent structures and was not run.

\section{Hosted inference and resource accounting}
\label{app:hosted}
Groq uses a chat-completions interface and Gemini a native interface. Both obtain credentials from environment variables, not stored experiment settings or command-line arguments. The following measurements describe the experimental period, not current quotas, prices, or availability.

\textbf{Groq.} The recorded allowance for \texttt{llama-3.3-70b-versatile} was 100,000 tokens per day. Full-vocabulary prompts exposed 29 amines, 74 aldehydes, and 14 isocyanides and consumed about 1,850 tokens per call, or 148,000 per 80-call seed before retries. A second independent key encountered the same ceiling. Sampling eight values per slot reduced consumption to about 813 tokens per call and 65,000 per seed. Output validation still used the complete vocabulary, preserving validity but changing each call's visible choices.

The allowance recovered over an observed rolling 24-hour window. Three daily attempts failed before a fourth succeeded; large quota probes were replaced by trivial probes to conserve tokens. Three random-split seeds accumulated over several days, all with zero rejection. Two underperformed both baselines; one exceeded GA only. A later catalog query reportedly omitted the model, but no dated endpoint snapshot supports a present-availability claim. An eight-seed, three-split sweep was not completed.

\textbf{Gemini.} The recorded \texttt{gemini-flash-latest} allowance was 20 requests per day in its free-tier per-project, per-model bucket. The separate \texttt{gemini-flash-lite-latest} bucket enabled eight seeds on all three splits. These completed results supersede partial runs. Flash-Lite is not the flagship model or an established frontier-capability benchmark.

\textbf{Rate limiting.} Hosted calls are paced about seven seconds apart. HTTP 429 responses invoke Groq's \texttt{Retry-After} header or an empirical Gemini delay, with a configured retry limit of eight before raising an error rather than silently replacing the model. Pacing controls bursts, not daily token totals; compression and quota recovery address the latter. Nominal workload costs exclude multiday scheduling and some retries.

\textbf{Domain-adapted models.} TxGemma was not evaluated because of access gating and insufficient 16\,GB VRAM for its 27B chat model. Prebuilt ChemLLM and ChemDFM checkpoints were unavailable in the local backend. Meditron 7B, continued-pretrained on biomedical and clinical text from Llama 2, repeatedly returned unrelated tool-shaped JSON, invalid reagent structures, or literal vocabulary placeholders. Its 100\% rejection and fallback rate precluded a meaningful sweep. Per-run rejection logging distinguishes this failure from genuine model performance.

\section{Diagnostic example and experiment provenance}
\label{app:example}
\subsection{Repeated proposals within one pool}
A default-prompt Llama 8B pool of eight candidates, initialized from the recorded six-fragment random archive, repeats the first triple six times and each remaining triple once. All are labeled crossover. Ellipses in the last two isocyanide slots are retained abbreviations, not reconstructed molecules.
\begin{Verbatim}[fontsize=\small,breaklines,breakanywhere]
A=CN(C)CCN                         crossover x6
B=O=C(/C=C\CCCCCCC)OCCCCCC=O
C=[C-]#[N+]CCCCCCCC/C=C\CCCCCCCC
A=NN1CCCCC1                       crossover x1
B=C#CCCCCCCCCC(=O)OCCC=O
C=...
A=NC1CN2CCC1CC2                    crossover x1
B=CCCCC/C=C\C/C=C\CCCCCCCC(=O)OCCC=O
C=...
\end{Verbatim}
Within-pool repetition differs from archive duplication. Archive-only de-duplication removes existing Pareto matches without backfill but does not ensure mutual uniqueness among novel proposals.

\subsection{Scope and reproducibility}
The experiment record distinguishes oracle-validation outputs, per-seed search results, acquisition-disabled diagnostics, diversity prompting, and the corrected acquisition dispatch. These provenance records describe the study; the code that produced them is publicly released (Code availability statement), though per-seed raw query traces beyond the aggregated statistics reported here are available from the corresponding author upon request rather than bundled in the repository. Initial Llama and proposal-diversity diagnostics used four seeds before the eight-seed, three-split sweep. Later panel aggregates summarize larger samples and need not equal pilot rates. Acquisition-disabled results are diagnostic, not a complete eight-seed factorial study.

Budget regimes are sequential controls, not interchangeable descriptions: $B=400$ tests saturation on the original 1,200-candidate space; $B=80$ reduces the budget before and after expansion to 30,044 candidates; $B=40$ defines the primary LLM protocol; and $B=15$ defines the initial four-seed de-duplication pilot. Primary inference uses eight random-split comparisons. Other tests are secondary or exploratory, including eight-seed de-duplication.

The recorded reproducibility materials include exact partition indices with split-generation seed zero and 960 training and 240 test molecules per split, plus a manifest mapping tables to experiment settings and outputs. GP settings are a Tanimoto kernel, diagonal regularization $\alpha=10^{-6}$, internal target normalization, and random seed zero. Bootstrap intervals use 10,000 resamples, 95\% two-sided percentile bounds, and random seed zero. Objective definitions and the single-assay exception appear in Section~\ref{sec:objectives}. The internal null-result interpretation was not independently timestamped. The code release (Code availability statement) is not anonymized, consistent with this being a non-anonymous journal submission.

\subsection{Eight-seed de-duplication and exploratory EHVI}
\label{app:dedup-full8}
For the random-split $B=40$ archive-de-duplication intervention, Llama 8B hypervolumes are 135.463, 90.001, 134.490, 123.479, 100.779, 100.779, 167.939, and 100.779. Matched GA values are 159.384, 109.838, 135.463, 114.082, 131.714, 201.687, 134.490, and 152.093; random-search values are 139.085, 155.060, 166.663, 162.617, 124.540, 175.418, 130.415, and 155.060.

The independently replicated no-de-duplication control gives Llama values 100.779, 80.137, 100.779, 123.479, 100.779, 100.779, 100.779, and 128.397. Its GA values are 159.384, 112.501, 135.463, 100.518, 125.852, 199.167, 134.490, and 144.025; the random-search vector is unchanged. De-duplication applies to every method, so GA's nonzero duplication changes its trajectory, whereas no observed random-search duplicates are removed. Using the distinct GA control reproduces the median $-33.0$ and $p=0.023$ in Table~\ref{tab:dedup}. A prior transcription incorrectly reused the de-duplicated GA vector for this control; the corrected vectors retain the reported summary statistics.

De-duplicated seeds 4 and 7 terminate after 31 and 14 queries because of repeated empty pools, not runtime errors. Thus the nominal $B=40$ cap is not exhausted in these two cases.

The four-seed EHVI run gives Llama values 100.779, 88.219, 100.779, and 123.479; GA values 123.479, 112.501, 135.463, and 112.282; and random-search values 139.347, 155.060, 132.080, and 161.350. No Wilcoxon result is reported: the recorded invocation failed on the observed tie pattern, and an exact two-sided four-seed test cannot attain $p<0.05$ even without ties. Small $n$ limits significance resolution, rather than making every signed-rank calculation intrinsically invalid.

\subsection{Cold-split de-duplication and secondary ablations}
\label{app:dedup-other-splits}
At $n=4$, scaffold-cold control differences are $-34.22$ against GA and $-39.96$ against random search; de-duplication gives $-5.09$ and $-13.37$. Chemotype-cold controls give $-45.29$ and $-75.06$, changing to $-34.98$ and $-70.16$. All seeds exhaust $B=40$. No signed-rank results are reported for these exploratory comparisons; only median differences and bootstrap intervals are used. Scaffold-cold narrows by factors of about 6 and 3, compared with 1.3 and 1.07 for chemotype-cold. This asymmetry remains mechanistically unresolved and must be interpreted alongside chemotype-cold's failed RAW264.7 validity gate.

The $B=80$, $m=40$, $q=8$ Ledger and binary-filter ablations cover chemotype-cold and random splits, respectively, at eight seeds, not a complete cross-product. Recorded seed-zero traces report eight Ledger-ablation rounds with six final-round queries and nine binary-filter rounds with four final-round queries. These counts alone cannot reconcile $B=80$ with at most eight selected queries per round: partial final rounds do not explain the discrepancy. Complete initialization and query accounting are required before these trace counts can substantiate budget exhaustion. This limitation does not alter the reported equality of ablation outcomes.

\subsection{Reporting corrections}
Llama has four unadjusted significant comparisons, not five; together with Qwen 7B this is seven of twelve. Qwen 14B's positive random-split point is against GA. Groq's 7.1\% duplication is below the small-model and GA rates but above Qwen 14B's aggregate 3.5\%. Temperature 1.0's 72.2\% duplication is not below the separate 67.8\% default diagnostic. Seven to 21 de-duplicated pilot rounds versus three control rounds imply a 2.3--7-fold increase. The three implementation errors concern split-oracle reuse, mutation mismatch, and acquisition dispatch. These distinctions correct interpretation without changing measured entries.

\section{Baseline proposal-pool specification}
\label{app:baselines}
All methods share archive and vocabulary access, $m=10$, $q=5$, the evaluation oracle, and the acquisition function. By default, no method removes within-pool duplicates. Archive-only filtering, when enabled, removes exact Pareto matches before scoring; all-evaluated filtering also removes previously evaluated dominated candidates and within-pool repetitions.

Each method constructs its proposal pool differently from a common archive and vocabulary. The LLM proposer draws up to eight sampled Pareto members per call, mutates by having the model choose one slot to change, and crosses over by having the model recombine two archive parents; duplicate handling is none by default, with the common optional filters available to every method. Static-GA draws two uniformly sampled archive members without replacement as parents, mutates with probability 0.5 by redrawing one random slot uniformly, and crosses over by copying each slot from a uniformly chosen parent; duplicate handling is likewise none by default with the same optional filters available. Random search draws independent uniform vocabulary samples with no parent source, mutation, or crossover operator, and no duplicate handling is needed since its measured archive-duplicate rate is zero. The Thompson bandit refits per-slot Beta posteriors from the archive each round, draws independently from those posteriors categorically (its analog of mutation), performs no parent recombination (no crossover), and applies no duplicate handling by default.

Invalid LLM outputs are replaced individually by the rule-based proposer and counted in the rejection rate. By contrast, no method backfills candidates removed by verification or novelty filtering to restore a fixed pool size. If fewer than $q$ survive, acquisition queries $\min(q,B_t,|P_t'|)$; subsequent rounds replenish proposals. GA falls back to a uniform triple only with fewer than two archive members, including the initial empty archive. All methods use the same UCB or EHVI selection, not method-specific shortcuts.

\section{LLM inference and prompt specification}
\label{app:llm-repro}
Instruction-tuned identifiers are \texttt{llama3.1:8b}, \texttt{qwen2.5:7b}, and \texttt{qwen2.5:14b} for local Ollama inference, \texttt{gemini-flash-lite-latest} for Gemini, and \texttt{llama-3.3-70b-versatile} for Groq. Local runs use the backend library's default GGUF quantization; hosted APIs expose no local quantization control. The recorded aliases do not establish immutable checkpoint revisions.

No system prompt is used. Calls are single-turn and stateless, with context reinserted each time. Default temperature is 0.7; exploratory checks use 0.2 and 1.0. Top-$p$, top-$k$, and repetition penalties retain backend defaults. Local calls request JSON formatting and draw a call seed from the proposer's seeded generator; reproducibility remains backend-version dependent. Hosted calls request JSON in the prompt without server-side schema enforcement and allow 300 output tokens. The default retry limit is two, increased to eight for hosted runs; quota handling is described in Appendix~\ref{app:hosted}.

Reagent strings must match vocabulary entries character-for-character. Invalid outputs invoke individual rule-based fallback and are tracked through the rejection rate, which remains at most a few percent except for Meditron's 100\%. Default prompts expose the full vocabulary; Groq alone samples eight values per slot. Each call shows up to eight uniformly sampled archive members. There is no additional truncation or score summary. Consequently, the default proposer sees identities and archive membership, not objective values or an explicit maximization instruction.

The default user-turn template is reproduced below, with archive members and vocabulary lists interpolated at runtime. Line wrapping is typographic only.
\begin{Verbatim}[fontsize=\small,breaklines,breakanywhere]
You are the Proposer Agent in an evolutionary lipid-design loop. Candidates are Ugi
three-component-reaction lipids, each a triple (A, B, C):
A = amine, one of: [...]
B = aldehyde, one of: [...]
C = isocyanide, one of: [...]

Current Pareto archive (sample):
- A=... B=... C=...
[up to 8 lines]

Propose exactly ONE new candidate via crossover (recombining A/B/C from two archive
parents) or mutation (changing exactly one of A/B/C from a single archive parent, or
a random triple if the archive is empty). The value of A, B, and C in your output
MUST be copied EXACTLY (character for character) from the lists above - do not
paraphrase or alter the SMILES strings in any way.

Respond with ONLY a single JSON object shaped like:
{"A": "<one of the A values>", "B": "<one of the B values>", "C": "<one of the C values>",
  "provenance": {"operator": "crossover" or "mutation", "parents": ["..."]}}
\end{Verbatim}
The diversity variant adds a paragraph describing repetition and listing up to 20 evaluated triples to avoid, without other changes. The scored variant supplies sampled members' actual HeLa and RAW264.7 values and explicitly instructs maximization of both objectives.

\section{Equivalence and multiplicity}
\label{app:equivalence}
A non-significant difference test does not establish equivalence. The post hoc practical margin is 10\% of the baseline's median hypervolume on the relevant split. Table~\ref{tab:tost} reports the six available two-one-sided-test (TOST) results; it is not a complete equivalence analysis of every primary comparison.

\begin{table}[htbp]
\caption{Post hoc equivalence tests with a margin of 10\% of baseline median hypervolume. No reported test reaches $p<0.05$.}
\label{tab:tost}
\centering
\begin{tabular}{@{}llccc@{}}
\toprule
Model & Baseline & Threshold ($\pm$) & Mean diff & TOST $p$ \\
\midrule
Llama 3.1 8B & GA (random) & 15.3 & $-45.4$ & 0.947 \\
Llama 3.1 8B & RS (random) & 15.5 & $-39.2$ & 0.988 \\
Qwen2.5 14B & GA (random) & 15.3 & $-9.6$ & 0.354 \\
Qwen2.5 14B & RS (random) & 15.5 & $-3.4$ & 0.102 \\
Qwen2.5 14B & GA (scaffold-cold) & 13.7 & $+20.6$ & 0.630 \\
Qwen2.5 14B & RS (scaffold-cold) & 13.7 & $+25.8$ & 0.710 \\
\bottomrule
\end{tabular}
\end{table}

Qwen 14B therefore has neither evidence of a difference nor evidence of equivalence on random and scaffold-cold splits at this margin and $n=8$. Larger samples or a different justified margin would be needed to resolve practical equivalence. The margin was not prespecified before experimentation.

The primary multiplicity family contains the four fully replicated models against two baselines on the random split. Table~\ref{tab:holm} retains the raw $p$-values and adds valid Holm-adjusted values at family-wise $\alpha=0.05$.

\begin{table}[htbp]
\caption{Holm correction for the eight primary random-split comparisons. GA denotes static-GA and RS random search. The step-down procedure stops at the first failed threshold; tied minimum $p$-values cannot yield an isolated rejection.}
\label{tab:holm}
\centering
\begin{tabular}{@{}lccc@{}}
\toprule
Comparison & Raw $p$ & Ordered threshold & Adjusted $p$ \\
\midrule
Llama 8B vs.\ RS & 0.0078125 & 0.00625 & 0.0625 \\
Gemini vs.\ GA & 0.0078125 & 0.00714 & 0.0625 \\
Gemini vs.\ RS & 0.0078125 & 0.00833 & 0.0625 \\
Qwen 7B vs.\ RS & 0.015625 & 0.01000 & 0.078125 \\
Llama 8B vs.\ GA & 0.0546875 & 0.01250 & 0.21875 \\
Qwen 7B vs.\ GA & 0.078125 & 0.01667 & 0.234375 \\
Qwen 14B vs.\ GA & 0.7421875 & 0.02500 & 1.0000 \\
Qwen 14B vs.\ RS & 0.7421875 & 0.05000 & 1.0000 \\
\bottomrule
\end{tabular}
\end{table}

The minimum raw $p=0.0078125$ exceeds the first threshold, 0.00625, so no null hypothesis is rejected. A previous rank-by-rank interpretation incorrectly retained the third tied comparison while ignoring the step-down stopping rule. Correcting that interpretation weakens model-specific inferiority claims but leaves the conclusion of no demonstrated improvement unchanged. Displayed unadjusted tests elsewhere must not be read as family-wise evidence.

\section{Normalized hypervolume}
\label{app:normalized-hv}
Full-space hypervolume is obtained by embedding and evaluating all 30,044 candidates with each frozen oracle, without LLM calls or a search budget. It is a diagnostic denominator, not information exposed to a budgeted search. With an empty initial archive and initial hypervolume zero,
\begin{equation}
\mathrm{NormalizedHV}=\frac{\mathrm{HV}_{\mathrm{method}}}{\mathrm{HV}_{\mathrm{full\text{-}space}}}.
\end{equation}
\begin{table}[htbp]
\caption{Attainable hypervolume over the complete candidate space under each split-specific frozen oracle.}
\label{tab:full-space-hv}
\centering
\begin{tabular}{@{}lc@{}}
\toprule
Split & Full-space hypervolume \\
\midrule
Random & 258.2 \\
Scaffold-cold & 235.9 \\
Chemotype-cold & 199.5 \\
\bottomrule
\end{tabular}
\end{table}

In the primary random-split runs, Llama 8B reaches 106.4--129.5, or 41--50\% of the full-space value; GA reaches 103.8--219.7, or 40--85\%; random search reaches 124.5--175.4, or 48--68\%. These ranges refer to the primary runs, not intervention replicates. They do not establish a universal half-frontier upper bound for all models or interventions. At $B=40$, the query cap covers about 0.13\% of the space. Multiplying every paired difference on a split by the same positive normalization factor changes its scale, not the corresponding significance result.

\section{All-evaluated de-duplication}
\label{app:dedup-all-evaluated}
For random-split seeds 0--3 with $B=40$, stricter de-duplication gives Llama 8B values 135.463, 232.263, 134.490, and 135.463; GA values 182.084, 112.501, 135.463, and 94.317; and random-search values 139.085, 155.060, 166.663, and 162.617. The matched control gives Llama values 100.779, 80.137, 100.779, and 123.479; GA values 159.384, 112.501, 135.463, and 100.518; and the same random-search vector. The unchanged random-search outcomes apply to this matched comparison, not every separate experiment on the split. Seed 1's value of 232.263 is retained and strongly influences the exploratory result.

\section{Objective-scored prompting}
\label{app:v3scored}
The four scored-prompt hypervolumes are 96.167, 232.263, 135.463, and 128.397, at $B=40$ on the random split. GA values are 159.384, 112.501, 135.463, and 100.518; random-search values are 139.085, 155.060, 166.663, and 162.617. These are the no-de-duplication baseline vectors, not the stricter-filter GA vector. The default-prompt Llama values are 100.779, 80.137, 100.779, and 123.479.

Seed 1 reaches 232.262563068695 in both scored prompting and all-evaluated de-duplication. The recorded per-round reagent sequences differ, excluding identical cached query trajectories. Convergence from a shared seed makes these two exploratory checks dependent; it does not justify treating their favorable medians as independent corroboration.

\section{HeLa-only chemotype-cold reanalysis}
\label{app:helaonly}
The single-objective restriction uses the primary $n=8$, $B=40$, $m=10$, $q=5$ protocol, the HeLa readout, and reference point $(-5.0)$. Llama 8B hypervolumes are 10.994, 10.609, 10.527, 10.905, 8.826, 11.422, 12.917, and 11.422. Both baselines give 13.878, 14.147, 17.542, 17.159, 13.142, 16.706, 12.766, and 14.959. Their identical vectors reflect equal best-found values under this metric. Llama's paired median difference is $-3.93$ against each, with unadjusted $p=0.015625$. The implementation explicitly makes the objective heads configurable for this exception rather than retaining the historical two-head default.

\section{Evidence-use checks}
\label{app:leakage8b}
\subsection{Leakage probe (Llama 3.1 8B)}
Llama 3.1 8B zero-shot scores the same 50 held-out random-split lipids used for the Qwen3-8B control, from 1 to 10, with no labels. At temperature zero, rejection is zero but all 50 emitted scores equal 8.0 -- a degenerate constant output for which Spearman correlation against the fitted oracle's potency prediction (itself $\rho=0.683$ on this split, Section~\ref{sec:validity}) is undefined rather than computably near zero; this is tracked explicitly (\texttt{degenerate\_constant\_output=true} in the recorded result) rather than silently reported as a correlation of zero. Repeating the identical protocol at temperature 0.7 (same 50 molecules, zero rejection) breaks the degeneracy: three distinct scores are emitted (6, 7, 8), and Spearman correlation is $\rho=-0.0065$, $p=0.96$. This confirms the temperature-zero degeneracy was a greedy-decoding artifact of this specific prompt and model, not evidence that the question cannot be meaningfully answered, and the non-degenerate result is consistent with the separate Qwen3-8B zero-shot control ($\rho=0.014$, $p=0.92$, $n=50$): both models show a measured near-zero correlation once decoding is non-degenerate, supporting the conclusion that no exploitable pretrained potency signal is recoverable from either model under zero-shot scoring.

\subsection{Rank-only prompt (\texttt{v4\_rank})}
\label{app:v4rank}
The prompt lists up to eight sampled archive members in rank order by summed potency and RAW264.7 score, with no numeric magnitudes shown, and otherwise matches the default template. On four matched random-split seeds at $B=40$, hypervolumes are 128.397, 227.271, 142.046, and 128.397, against the GA baseline vector 159.384, 112.501, 135.463, and 100.518, and random-search vector 139.085, 155.060, 166.663, and 162.617, giving per-seed differences against GA of $-31.0$, $+114.8$, $+6.6$, and $+27.9$ (median $+17.2$) and against random search of $-10.7$, $+72.2$, $-24.6$, and $-34.2$ (median $-17.7$); both CIs cross zero at $n=4$.

\subsection{Chain-of-thought rationale (\texttt{v5\_cot})}
\label{app:cot}
The schema requires a non-empty \texttt{"rationale"} field of one to two sentences describing the structural reasoning behind the proposed crossover or mutation; outputs without one are rejected and retried, counted toward the rejection rate. On four matched random-split seeds at $B=40$, hypervolumes are 128.397, 106.697, 128.397, and 104.358, giving median differences $-6.4$ against GA and $-43.3$ against random search. Rejection remains under 5\%, so the requirement is enforced rather than decorative.

\subsection{Score-shuffle perturbation (\texttt{v6\_shuffled})}
\label{app:shuffle}
The prompt displays real potency and RAW264.7 values exactly as in the scored prompt, but the values are independently shuffled across the displayed fragments before rendering, breaking the true value-to-fragment correspondence while preserving the set of displayed numbers. On four matched random-split seeds at $B=40$, hypervolumes are 96.876, 232.263, 135.463, and 128.397, giving median differences $+13.9$ against GA and $-32.7$ against random search -- matching the correctly scored prompt's $+13.9$ and $-32.7$ (Table~\ref{tab:scored-prompt}) to one decimal place. Direct per-query comparison of reagent triples against the correctly scored run's \texttt{queries.csv} finds real divergence in three of four seeds: seed 0 differs on 28/40 queries, seed 1 on 8/40, seed 2 on 0/40 (fully identical trajectory), and seed 3 on 7/40. The attainable Pareto frontier under this budget and split is narrow enough that distinct query trajectories, including trajectories built from deliberately false score mappings, often converge on the same final hypervolume. Aggregate equivalence therefore reflects frontier concentration, not proof that scores go unused.

\subsection{Evidence-conflict probe}
\label{app:conflict}
A fixed eight-member sample is drawn from the random-split training partition's evaluated archive, independently for each of three probe-sampling seeds (0, 1, 2). Within each sample, true rank is computed by summed potency and RAW264.7 score; displayed rank labels are the true ranking inverted (true-best shown as rank 8, true-worst shown as rank 1). This falsified ranking is rendered in the \texttt{v4\_rank} template and issued to Llama 3.1 8B 40 times independently per sample (temperature 0.7, distinct call seeds), each call producing one proposal via the proposer's own parsing path, not the full search loop. Each proposal's nearest neighbor among the eight shown members is found by minimum fragment-slot edit distance; calls that tie between multiple equally near members or resolve to no identifiable nearest member are excluded from the per-sample tally.

Table~\ref{tab:conflict-samples} reports all three samples. The three samples do not agree: the fraction of resolved calls nearest the falsely labeled-best member (the true worst) is 14.8\%, 84.2\%, and 46.2\%, and the fraction of calls that tie or fail to resolve to a unique nearest member is also highly variable (13, 2, and 27 of 40) -- plausibly reflecting how similar the eight randomly drawn archive members happen to be to each other in a given sample, which is not controlled across draws. The one pattern that holds in all three samples is that the genuinely best member is almost never the nearest neighbor regardless of its (false) label: 0, 3, and 0 of resolved calls, pooled 3/78 (3.8\%), well below the 12.5\% uniform base rate for an 8-member set. A single fixed-archive run of this probe therefore cannot support a specific claim about what proposals are tracking when the true and false quality signals conflict -- the evidence here rules out "tracks true quality despite the false label" but does not establish a replicated alternative; seed 1's 84.2\% false-label-following, taken alone, would have suggested near-complete capture by the false label, while seed 0's 14.8\% would have suggested near-independence from either label, illustrating why this probe needs more than one fixed sample before its outcome is reported as a stable finding.

\begin{table}[htbp]
\caption{Evidence-conflict probe across three independently drawn fixed archive samples, 40 calls each. ``Other'' denotes a resolved nearest neighbor that is neither the true-best nor the falsely labeled-best (true-worst) member.}
\label{tab:conflict-samples}
\centering
\begin{tabular}{@{}lccccc@{}}
\toprule
Sample seed & Resolved/40 & Nearest true-best & Nearest false-best & Nearest other & Tie/unresolved \\
\midrule
0 & 27 & 0 (0\%) & 4 (14.8\%) & 23 (85.2\%) & 13 \\
1 & 38 & 3 (7.9\%) & 32 (84.2\%) & 3 (7.9\%) & 2 \\
2 & 13 & 0 (0\%) & 6 (46.2\%) & 7 (53.8\%) & 27 \\
\midrule
Pooled & 78 & 3 (3.8\%) & 42 (53.8\%) & 33 (42.3\%) & 42 \\
\bottomrule
\end{tabular}
\end{table}

The pooled row is reported for completeness but should not be read as a single stable rate: it averages over three samples whose individual false-label-following rates span a 69-percentage-point range, so the pooled 53.8\% is an artifact of seed 1's outlier sample rather than a representative summary. The pooled true-best rate (3.8\%) is the only quantity that is consistent in direction across all three samples and is therefore the one load-bearing conclusion from this probe.

\subsection{Genuine reasoning-mode proposer}
\label{app:reasoning}
Ollama's native \texttt{think=true} field is enabled for \texttt{qwen3:8b}, a mechanism distinct from the prompted \texttt{v5\_cot} rationale; requesting JSON-formatted output simultaneously with \texttt{think=true} was found to silently suppress all reasoning output and is therefore not used, with JSON recovered instead by a hardened parser tolerant of markdown-fenced or prose-wrapped output. Under the default (unscored) prompt, four matched random-split seeds at $B=40$ give Llama-comparable baselines (GA: 159.384, 112.501, 135.463, 100.518; random search: 139.085, 155.060, 166.663, 162.617) against reasoning-mode hypervolumes 101.998, 136.623, 93.679, and 89.201, giving median differences $-26.55$ (CI $[-57.39,24.12]$) against GA and $-55.04$ (CI $[-73.42,-18.44]$) against random search. Diagnostics over 160 calls show 155 produced a reasoning trace (mean length 5,669 characters) and 0 mentioned potency or the objective. Under the scored prompt, all four seeds completed after raising Ollama's per-call timeout to 300s (from the 120s default), following one timeout-related failure on the first attempt. Scored-prompt hypervolumes are 101.203, 98.645, 111.197, and 111.197 against the same GA and random-search baseline vectors as above, giving median differences $-19.1$ (GA) and $-53.4$ (random search), both CIs crossing zero. Diagnostics over 160 calls show 152 produced a reasoning trace (mean length 6,877 characters) and 124 (77.5\%) mentioned potency, versus 0/155 under the default prompt -- a sharp rise in measured engagement with the displayed evidence that does not translate into a hypervolume advantage over non-reasoning Llama 8B under the identical scored prompt (Table~\ref{tab:scored-prompt}). Per-seed wall time for the default-prompt condition ranges from 2,801 to 5,290 seconds, roughly 5--10$\times$ the non-reasoning Llama or Qwen latency in Table~\ref{tab:original-6}.

\subsection{Prompt-variant latency}
\label{app:variant-latency}
A dedicated timing run issues eight fresh Llama 3.1 8B calls per prompt variant against a fixed eight-member synthetic archive, measuring wall-clock seconds per call under identical local GPU conditions within one session. Mean seconds per call are 11.85 (default), 10.59 (\texttt{v3\_scored}), 9.75 (\texttt{v4\_rank}), 15.60 (\texttt{v5\_cot}), and 13.99 (\texttt{v6\_shuffled}). This fresh run's default-prompt mean (11.85s) is higher than the 4.0s reported in Table~\ref{tab:original-6}, reflecting different GPU load/warm-state conditions between sessions rather than a changed model or prompt; absolute values are therefore not directly comparable across the two measurements, but the relative ordering within this single session is informative. The rationale-requiring \texttt{v5\_cot} variant is the slowest, consistent with longer generated output and occasional rejection-retry overhead from the enforced non-empty rationale field; the shuffled and rank-only variants are within the same order of magnitude as the default and scored prompts, since none of these changes to the prompt content substantially change expected output length.

\section{Structural inspection: additional detail}
\label{app:structures}
\subsection{3D rendering pipeline}
Three-dimensional conformers for Figure~\ref{fig:structures} are generated from each best-found SMILES by RDKit's ETKDG embedding followed by MMFF94 force-field geometry optimization (2,000 iterations), using the main project Python environment, and written to PDB files. These are rendered as ball-and-stick images (PyMOL's \texttt{util.cbaw} coloring, white background, ray-traced) using open-source PyMOL installed via conda-forge in a separate, isolated conda environment created specifically for this rendering step. This isolation is deliberate: PyMOL's dependency chain requires a newer numpy than the numpy version pinned throughout the rest of this project's GP, active-learning, and statistical-testing code, and the rendering step does not need to share a Python environment with any numerical result reported elsewhere in this paper. No lipid geometry, charge state, or protonation beyond RDKit's default valence-based structure generation is assumed; this is a visualization of real queried reagent identities, not a conformational or energetic claim about the physical lipid.

\subsection{Best-found lipids by method}
For the objective-scored-prompt check's seed 1, random split, Table~\ref{tab:best-lipids} gives the best-found triple by frozen-oracle predicted HeLa+RAW264.7 sum for each method. Each is the arg-max over that method's 40 queried candidates on this seed only, not a claim about the global optimum or about any other seed; different seeds and methods can and do find different best candidates.

\begin{table}[htbp]
\caption{Best-found reagent triple per method, objective-scored-prompt check, random split, seed 1. SMILES strings wrap at arbitrary characters (\texttt{seqsplit}) since they contain no natural word breaks.}
\label{tab:best-lipids}
\centering
\small
\begin{tabularx}{\linewidth}{@{}lXXXc@{}}
\toprule
Method & A (amine) & B (aldehyde) & C (isocyanide) & Predicted (HeLa, RAW), sum \\
\midrule
Llama 8B (scored prompt) & \seqsplit{CCN(CC)CCCN} & \seqsplit{CCCCCC(C)CCC(=O)OCCCC=O} & \seqsplit{[C-]\#[N+]CCCCCCCC/C=C\textbackslash CCCCCCCC} & (12.0, 7.4), 19.3 \\
Static-GA & \seqsplit{NCCN1CCCCC1} & \seqsplit{CCCCCCCCCCCCCCCC=O} & \seqsplit{[C-]\#[N+]CCCCCCCCCCC} & (8.6, 2.9), 11.5 \\
Random search & \seqsplit{CCCCN(CCCC)CCCN} & \seqsplit{CC(CCC(OCCCCCC=O)=O)CCCCC} & \seqsplit{[C-]\#[N+]CCCCCCCCCCCCCCCC} & (9.7, 5.6), 15.2 \\
\bottomrule
\end{tabularx}
\end{table}

\subsection{Fragment-importance methodology and full range}
\label{app:fragment-importance}
For each of the 29 amine, 74 aldehyde, and 14 isocyanide values in the full extended vocabulary, every full-vocabulary candidate (30,044 total) containing that value is identified, and the frozen random-split $\mathcal O_{\mathrm{GT}}$'s predicted HeLa+RAW264.7 sum is averaged across those candidates. Values appearing in fewer than 20 candidates are excluded to avoid small-sample means (this excludes none of the 29 amine or 14 isocyanide values, and a subset of the 74 aldehyde values). This is computed directly from the oracle's posterior mean over the full combinatorial space -- no search is run, no budget is charged, and no counterfactual pair or significance test is constructed, unlike the Hypothesis Ledger's intended design (Appendix~\ref{app:agents}).

Overall mean (across all 30,044 candidates) is 7.02 with standard deviation 1.41. The amine slot shows the widest range (3.20--9.61, a 6.4-point spread, dominated by cyclic tertiary amines with short linkers at the top and more flexible or doubly cyclic diamine-type structures at the bottom). The aldehyde slot's range is 3.49--8.93 (5.4 points); the five lowest-scoring values are all doubly unsaturated (diene) long-chain tails, and the five highest are saturated, branched, shorter-chain linkers. The isocyanide slot's range is narrow, 6.83--7.38 (0.55 points), and is dominated by tail length rather than unsaturation, though the two highest-scoring values are both mono-unsaturated. This means the apparent unsaturation association at slot C is both weaker in magnitude and opposite in direction to the one at slot B, so it should not be read as evidence for a general unsaturation-potency rule in this chemistry; any such claim would require an independently validated structure-activity study, not a correlational read of one fitted oracle.

\section{Literature-lipid domain-coverage check}
\label{app:lit-lipids}
Three clinically deployed ionizable lipids are checked against this paper's real \texttt{decompose\_lipid} routine (validated to exactly round-trip all 1,200 labeled AGILE lipids, Section~\ref{sec:protocol}): MC3 (DLin-MC3-DMA, PubChem CID 49785164), SM-102 (PubChem CID 126697616), and ALC-0315 (PubChem CID 122666778). Canonical SMILES for each were obtained from their respective PubChem records and independently verified by computing each SMILES's molecular formula with RDKit and matching it exactly against the formula listed on that compound's own PubChem page (C43H79NO2, C44H87NO5, and C48H95NO5, respectively) before use, rather than trusting a search-engine summary.

All three SMILES parse validly in RDKit, but none match the Ugi-3CR core substructure the decomposition routine requires (\texttt{decomposes\_as\_ugi3cr=False} for all three, recorded in \texttt{results/e\_literature\_lipid\_check.json}). This is the expected outcome: MC3 is a dilinoleyl ester tertiary amine, SM-102 and ALC-0315 are branched-ester tertiary amines with hydroxyl/ether linkers -- none are products of an Ugi three-component reaction, so none should be expected to share this paper's core scaffold. The check is reported because it was previously named in the main text as a planned but unreported comparison (Recall@$k$ against MC3/SM-102/ALC-0315, Section~\ref{sec:protocol}); running it and disclosing the negative result removes that gap rather than leaving a stated plan silently unaddressed. No oracle potency prediction is computed for these three compounds, since the frozen $\mathcal O_{\mathrm{GT}}$ requires a Morgan-fingerprint embedding of an assembled Ugi-3CR fragment triple and these molecules supply no such triple.

\section{Panel interpretation}
Figure~\ref{fig:architecture} presents the budget-accounted architecture; Algorithm~\ref{alg:leaf} and Section~\ref{sec:objectives} define implemented budget accounting and endpoints. Hypothesis generation need not query the oracle, but counterfactual validation does. In the reported experiments, the Verifier is deterministic and the Ledger remains empty.

Figure~\ref{fig:performance} reports held-out validity, median paired differences with bootstrap intervals, the zero reference line, unadjusted significance markers, and proposal diversity. Minimum edit distance counts changed reagent slots, from zero to three. Random search's zero duplication is measured, not missing. Qwen 14B is the low-duplication exception among the four replicated models. Gemini's aggregate duplication exceeds GA's, but its chemotype-cold value of 19.2\% is below GA's reported 22.5\%; split-specific results take precedence over blanket comparisons. Pilot and aggregate diagnostics remain distinct.

Figure~\ref{fig:diagnostics} retains mean learning curves with standard deviations, early baseline gains and LLM flattening, log-scale nominal workload costs, and the non-significant 60-run duplication correlation. Its heatmaps include all 36 eight-seed comparisons, including the diversity prompt and non-LLM Thompson bandit. Groq's three-seed experiment is excluded from those heatmaps, not treated as a complete eight-seed result.

\section{Complete numerical tables}
\label{app:tables}
These eight tables retain every numerical entry in the original manuscript's tables. Negative paired hypervolume differences favor the baseline. All tests are two-sided and unadjusted; eight seeds allow significance but do not guarantee statistical power. ``GA'' denotes static-GA and ``RS'' denotes random search.

\begin{table}[htbp]
\caption{Llama 3.1 8B: paired hypervolume differences and Wilcoxon tests ($n=8$).}
\label{tab:original-1}
\centering
\begin{tabular}{@{}lccc@{}}
\toprule
Comparison & Random & Scaffold-cold & Chemotype-cold \\
\midrule
LLM vs.\ GA & $p=0.055$ & $p=0.078$ & $p=0.039$ \\
Median & $-47.1$ & $-24.8$ & $-31.2$ \\
LLM vs.\ RS & $p=0.0078$ & $p=0.039$ & $p=0.016$ \\
Median & $-42.0$ & $-28.8$ & $-46.3$ \\
\bottomrule
\end{tabular}
\end{table}

\begin{table}[htbp]
\caption{Qwen2.5 7B: paired hypervolume differences and Wilcoxon tests ($n=8$).}
\label{tab:original-2}
\centering
\begin{tabular}{@{}lccc@{}}
\toprule
Comparison & Random & Scaffold-cold & Chemotype-cold \\
\midrule
Qwen vs.\ GA & $p=0.078$ & $p=0.109$ & $p=0.016$ \\
Median & $-40.9$ & $-21.6$ & $-29.2$ \\
Qwen vs.\ RS & $p=0.016$ & $p=0.195$ & $p=0.023$ \\
Median & $-29.1$ & $-34.3$ & $-43.4$ \\
\bottomrule
\end{tabular}
\end{table}

\begin{table}[htbp]
\caption{Qwen2.5 14B: paired hypervolume differences and Wilcoxon tests ($n=8$).}
\label{tab:original-3}
\centering
\begin{tabular}{@{}lccc@{}}
\toprule
Comparison & Random & Scaffold-cold & Chemotype-cold \\
\midrule
Qwen vs.\ GA & $p=0.742$ & $p=0.313$ & $p=0.0078$ \\
Median & $+1.8$ & $+18.3$ & $-55.1$ \\
Qwen vs.\ RS & $p=0.742$ & $p=0.250$ & $p=0.0078$ \\
Median & $-4.3$ & $+33.9$ & $-77.0$ \\
\bottomrule
\end{tabular}
\end{table}

\begin{table}[htbp]
\caption{Initial proposal diagnostics (four seeds, one split). These are pool-proposal statistics, not directly measured fractions of selected oracle queries.}
\label{tab:original-4}
\centering
\begin{tabularx}{\linewidth}{@{}Xccc@{}}
\toprule
Method & Crossover label & Archive duplicate & Edit distance \\
\midrule
Llama 3.1 8B & 99.7\% & 67.8\% & 0.587 \\
Qwen2.5 7B & 56.2\% & 62.5\% & 0.650 \\
Static-GA & 0\%$^\dagger$ & 22.5\% & 1.516 \\
Random search & 0\%$^\dagger$ & 0\% & 2.812 \\
\bottomrule
\end{tabularx}
\par\footnotesize $^\dagger$The non-LLM methods use different internal operator labels (e.g., \texttt{ga\_crossover\_mutation}); zero does not mean zero recombination.
\end{table}

\begin{table}[htbp]
\caption{Matched de-duplication check: Llama 3.1 8B, random split, $B=15$, $n=4$. All 95\% CIs cross zero.}
\label{tab:original-5}
\centering
\begin{tabular}{@{}lcc@{}}
\toprule
Comparison & No de-duplication & With de-duplication \\
\midrule
LLM vs.\ static-GA & Median $-17.6$ & Median $-4.7$ \\
LLM vs.\ random search & Median $-15.1$ & Median $+2.3$ \\
\bottomrule
\end{tabular}
\end{table}

\begin{table}[htbp]
\caption{Measured candidate latency and 80-call workload cost. Times reproduce the original table; hosted retries, quota delays, and multi-day accumulation are additional.}
\label{tab:original-6}
\centering
\small
\begin{tabularx}{\linewidth}{@{}X p{0.17\linewidth} p{0.22\linewidth} p{0.17\linewidth}@{}}
\toprule
Method & Backend & Seconds per candidate & 80-call workload \\
\midrule
Llama 3.1 8B & Local GPU, Ollama & $\approx4.0$ & $\approx5.3$ min \\
Qwen2.5 7B & Local GPU, Ollama & $\approx10.3$ & $\approx13.7$ min \\
Qwen2.5 14B & Local GPU, Ollama & $\approx10.3$ & $\approx13.7$ min \\
Gemini Flash-Lite & Hosted, free tier & $\approx5.1$ & $\approx6.8$ min \\
Llama 3.3 70B & Hosted, free tier & $\approx0.6$ raw; $\approx7$ paced$^\dagger$ & $\approx9.3$ min$^\dagger$ \\
Static-GA / random search & None (GP only) & $<0.01$ & $<1$ s \\
\bottomrule
\end{tabularx}
\par\footnotesize $^\dagger$Pacing addresses per-minute throughput; a separate 100,000-token daily allowance limits total runs. Pacing alone does not reduce the tokens needed per run. The non-LLM row denotes the matched search workload, not 80 LLM calls.
\end{table}

\begin{table}[htbp]
\caption{Thompson-sampling bandit: paired hypervolume differences and Wilcoxon tests ($n=8$).}
\label{tab:original-7}
\centering
\begin{tabular}{@{}lccc@{}}
\toprule
Comparison & Random & Scaffold-cold & Chemotype-cold \\
\midrule
Bandit vs.\ GA & $p=0.742$ & $p=0.945$ & $p=0.028$ \\
Median & $-16.2$ & $+5.9$ & $-18.0$ \\
Bandit vs.\ RS & $p=0.844$ & $p=1.000$ & $p=0.039$ \\
Median & $+0.9$ & $+3.4$ & $-35.2$ \\
\bottomrule
\end{tabular}
\end{table}

\begin{table}[htbp]
\caption{Gemini Flash-Lite: paired hypervolume differences and Wilcoxon tests ($n=8$).}
\label{tab:original-8}
\centering
\begin{tabular}{@{}lccc@{}}
\toprule
Comparison & Random & Scaffold-cold & Chemotype-cold \\
\midrule
Gemini vs.\ GA & $p=0.0078$ & $p=0.0078$ & $p=0.0078$ \\
Median & $-91.8$ & $-66.6$ & $-24.8$ \\
Gemini vs.\ RS & $p=0.0078$ & $p=0.0078$ & $p=0.0078$ \\
Median & $-85.5$ & $-71.1$ & $-60.2$ \\
\bottomrule
\end{tabular}
\end{table}


\end{document}
